\providecommand{\toplevelprefix}{../..}  %
\documentclass[../../book-main.tex]{subfiles}

\begin{document}

\chapter{Aprendizaje de representaciones mediante compresión con pérdidas}
\label{ch:lossy-compression}

\begin{quote}
    \hfill ``\textit{Comprimimos para aprender y aprendemos para comprimir}.''

    $~$\hfill -- \textit{High-Dimensional Data Analysis}, Wright y Ma, 2022
\end{quote}
\vspace{5mm}

Recapitulemos lo que hemos visto hasta ahora. Hemos mostrado que es
posible identificar y acceder a una distribución con un soporte de
dimensiones arbitrariamente bajas, como la que se ilustra en la
\Cref{fig:representation-learning} (izquierda), transformando gradualmente
la distribución de ruido puro (alta entropía) en la distribución que 
se desea aprender mediante \textit{eliminación de ruido iterativo}. Hemos
mostrado que el filtro de eliminación de ruido en cada iteración está directamente relacionado con
el gradiente de la densidad logarítmica de la distribución mediante la fórmula de Tweedie
\eqref{eq:tweedie}. Demostramos que, en la práctica, ese filtro de eliminación de ruido
\(\bar{\vx}_{\theta}\) puede estimarse utilizando un número finito de muestras. Por lo tanto,
desarrollamos un método general para aprender o aproximarse a una distribución de datos.

No obstante, en esta metodología, la
codificación de la distribución está implícita en la forma funcional
y los parámetros del filtro de ruido, si los hubiera. De hecho, los lectores más atentos habrán
notado que, para una distribución general, nunca hemos
especificado explícitamente cuál debería ser la forma funcional del filtro de ruido. En la práctica,
la gente suele modelarlo con alguna red neuronal profunda que tiene una
arquitectura empíricamente diseñada. Además, aunque sabemos que el
proceso de eliminación de ruido anterior reduce la entropía, no sabemos en qué
medida, ni conocemos la entropía de las distribuciones intermedias y finales.
Por lo tanto, sería muy útil conocer la forma (óptima)
de la operación de eliminación de ruido y cómo medir el cambio en
la entropía. Por último, al final del capítulo anterior demostramos que si
el proceso de eliminación de ruido reduce la entropía \textit{en exceso}, esto conduce
a que solo se aprendan los datos de entrenamiento, sin posibilidad de interpolación o
extrapolación. La innovación metodológica que resuelve todas
estas cuestiones se desarrollará en el próximo capítulo.

\paragraph{Medida de información.}
Recordemos que nuestro objetivo general es modelar datos procedentes de una distribución continua
con un soporte de baja dimensión. Si queremos identificar el
modelo más simple que genera los datos, podríamos considerar tres
medidas típicas de parsimonia: la dimensión, el volumen o la
entropía diferencial. Si usamos la dimensión, es obvio que el mejor modelo para un
conjunto de datos dado es la propia distribución empírica, que es
de dimensión cero.\footnote{En la Introducción
    \Cref{ch:intro}, sostuvimos que un objetivo importante de la inteligencia es
    buscar estructuras de baja dimensión en los datos observados. Sin embargo,
    desde el punto de vista técnico, no se debe interpretar el término
    «de baja dimensión» de manera literal. El objetivo de este capítulo es
    aclarar con precisión qué estructuras de baja dimensión debemos
buscar y cómo hacerlo.} Si, en cambio, utilizamos la entropía de los datos, entonces
para todas las distribuciones con soporte de baja dimensión, la entropía diferencial
es siempre infinito negativo. O bien, si utilizamos el volumen del
espacio que abarcan los datos, el volumen del soporte de cualquier
modelo de baja dimensión es siempre cero. Entonces, entre todas las distribuciones con
un soporte de baja dimensión que podrían haber generado las mismas
muestras de datos, ¿cómo podemos decidir cuál es la mejor basándonos en estas
medidas de parsimonia que no permiten distinguirlas entre sí en absoluto?

En este capítulo, analizamos un marco que mitiga la dificultad técnica mencionada anteriormente
al asociar la distribución aprendida con un
esquema explícito y computable de codificación y decodificación, siguiendo y
haciendo más explícito el segundo enfoque sugerido al final de la
\Cref{sub:min_entropy}. Como veremos, este enfoque nos permite, en esencia,
aproximar con precisión la entropía de las distribuciones aprendidas
en términos de una \textit{tasa de codificación con pérdida} asociada
al esquema de codificación, tal como se describe y justifica (aunque no se formaliza)
al final de la \Cref{sec:diffusion_memorization}. En términos generales, esta medida 
puede considerarse como una noción generalizada
del volumen de la distribución de los datos empíricos. Con una medida de este tipo,
podemos cuantificar con precisión en qué medida se reduce la entropía
mediante una transformación de la distribución, por ejemplo, mediante un filtro de eliminación de ruido. Además,
con una medida de este tipo, podemos obtener, al menos en principio,
el operador óptimo de eliminación de ruido que reduce la entropía más rápidamente.

\paragraph{El aprendizaje de representaciones.}
Aunque el proceso de difusión y eliminación de ruido presentado en el
capítulo anterior nos permite, en principio, aprender una distribución
arbitraria de baja dimensión, no ofrece una representación explícita
de dicha distribución. Como veremos, se accederá a la distribución
aprendida en repetidas ocasiones y se utilizará para tareas posteriores
(por ejemplo, como distribución a priori para la inferencia bayesiana en el
\Cref{ch:conditional-inference}). Por lo tanto, es conveniente que podamos
transformar la distribución de los datos, por ejemplo mediante una función $f(\vx,\theta)$, a
alguna forma estructurada estándar que facilite un acceso eficiente, como
una mezcla de subespacios lineales (incoherentes) o gaussianas de bajo rango, tal y como
se ilustra en la figura de la derecha de la \Cref{fig:representation-learning}. En
este capítulo formalizaremos este proceso como «aprendizaje de representaciones».

\begin{figure}
    \centering
    \includegraphics[width=0.8\linewidth]{chapters/chapter4/figs/manifold-linearization.pdf}
    \caption{Aprender una representación estructurada (derecha) de la distribución de los datos
        (izquierda). Obsérvese que, además de aprender (y
        muestrear) la distribución de los datos originales como hicimos en el
        anterior \Cref{ch:denoising}, aquí queremos aprender una transformación
        $f(\vx, \theta)$ para la distribución de datos, de modo que se obtenga una
        representación que tenga las propiedades geométricas o estadísticas deseadas,
    por ejemplo, lineal e incoherente.}\label{fig:representation-learning}
\end{figure}

Como veremos más adelante, para una distribución arbitraria, la medida de información anterior
(tasa de codificación) no se puede calcular (fácilmente). Sin embargo,
en el caso de una mezcla de subespacios o gaussianas de bajo rango, podemos obtener una
fórmula cerrada precisa para la medida. Al optimizar esta medida,
podremos encontrar la representación óptima. Como veremos en
el siguiente \Cref{ch:deep-networks}, la optimización de la medida también nos permite
derivar formas explícitas del operador que realiza la
transformación $f(\vx,\theta)$ de la manera más eficiente. Esto
dará lugar a una explicación basada en principios sobre la arquitectura de las redes profundas,
así como a diseños de arquitectura profunda más eficientes.

\section{Compresión mediante codificación con pérdidas}
\label{sec:lossy_compression}

En esta sección, desarrollamos la \textit{tasa de codificación con pérdida}
como una medida computable de complejidad para distribuciones con 
un soporte de baja dimensión, un contexto en el que las medidas clásicas como
la entropía, la dimensión y el volumen no logran distinguir entre los modelos
candidatos. En primer lugar, justificamos la necesidad de la codificación con 
pérdidas en relación con estas patologías (\Cref{sub:rate-dist}); a continuación, 
demostramos que la función de tasa-distorsión relaciona esta medida de la teoría 
de la codificación con la geometría del soporte de datos
por medio del recubrimiento esférico (\Cref{thm:covering-number-rate-distortion}).
Esta conexión constituye el puente entre los enfoques basados en la teoría de la información y los geométricos
para el aprendizaje de distribuciones de baja dimensión. Por último, desarrollamos
una aproximación exacta en forma cerrada de la tasa de codificación con pérdida para las distribuciones gaussianas
(\Cref{subsec:lossy DR}) y un algoritmo eficaz para agrupar mezclas de
distribuciones gaussianas de baja dimensión (\Cref{sec:clustering-Gaussians}), que se
generalizará a distribuciones no lineales en la siguiente sección.

\subsection{La necesidad de la codificación con pérdidas}
\label{sub:rate-dist}

Ya hemos abordado en varias ocasiones una dificultad: si al final
aprendemos la distribución a partir de un número finito de muestras y nuestra
clase de funciones de filtros de ruido contiene suficientes funciones, ¿cómo
nos aseguramos de que estamos tomando muestras de la distribución \textit{verdadera} (con
soportes de baja dimensión), en lugar de cualquier otra distribución que pueda
generar ese número finito de muestras con alta probabilidad? Veamos
algunas de las dificultades conceptuales y técnicas con ejemplos concretos.

\begin{example}[Volumen, dimensión y entropía]\label{eg:measures-of-complexity}
    Para el ejemplo que se muestra en la parte superior de la \Cref{fig:1d-line}, supongamos
    que hemos tomado algunas muestras de una distribución uniforme en una recta
    (por ejemplo, en un plano bidimensional). El volumen de la línea o de los conjuntos de muestras es cero.
    Desde el punto de vista geométrico, la distribución empírica sobre el conjunto de muestras finito
    obtenido es aquella de {\em dimensión mínima} capaz de generar
    dicho conjunto de muestras finito.\footnote{Un conjunto de muestras discretas tiene
    dimensión cero, mientras que la línea de apoyo tiene dimensión uno.}
    Pero esto parece estar en contraste con otra medida más de
    la complejidad: la entropía. La entropía (diferencial) de la línea es
    el infinito negativo, pero la entropía (discreta) de este conjunto de muestras
    es finita y positiva. Así pues, parece que nos encontramos ante una situación paradójica
    según estos criterios habituales de parsimonia o
    complejidad: no son capaces de diferenciar adecuadamente entre (los modelos de)
    distribuciones con soportes de baja dimensión y algunos parecen
    diferenciarlos incluso de maneras totalmente opuestas.\footnote{Por supuesto, en sentido estricto, la
    entropía diferencial y la entropía discreta no son directamente comparables.}
\end{example}

\begin{example}[Densidad]\label{eg:density} Consideremos los dos conjuntos de
    puntos de datos muestreados que se muestran en la \Cref{fig:1d-line}. Geométricamente,
    son esencialmente iguales: cada conjunto consta de ocho puntos
    y cada punto ha aparecido con una frecuencia igual a 1/8. La
    única diferencia es que, en el segundo conjunto de datos, algunos puntos están
    lo suficientemente «cerca» como para que se considere que tienen una mayor densidad alrededor
    de su respectivo «grupo». ¿Cuál de las dos es más relevante para la
    distribución real que podría haber generado las muestras? ¿Cómo podemos
    resolver esta ambigüedad a la hora de interpretar este tipo de distribuciones (empíricas)?
    \begin{figure}[t]
        \centering
        \includegraphics[width=0.7\linewidth]{\toplevelprefix/chapters/chapter4/figs/one-dim-distribution.png}
        \caption{Ocho puntos situados en una línea con diferentes
            configuraciones geométricas; distribuidos uniformemente (arriba) o muy
        agrupados (abajo).}
        \label{fig:1d-line}
    \end{figure}
\end{example}

Existe otra dificultad técnica relacionada con
la creación de un esquema explícito de codificación y decodificación para un conjunto de datos.
Dado un conjunto de datos muestreados en $\vX = [\vx_1, \ldots, \vx_N]$, ¿cómo
se puede diseñar un esquema de codificación que sea implementable en máquinas con
recursos computacionales y de memoria finitos? Nótese que incluso para representar un número real
cualquiera se necesita un número infinito de dígitos o bits. Por lo tanto,
cabe preguntarse si la entropía de una distribución es una medida directa
de la complejidad de su esquema de codificación (óptimo). Analizamos
este tema con otro ejemplo sencillo.
\begin{example}[Precisión] \label{eg:two-inrational}
    Consideremos una distribución discreta $\vX = [e, \pi]$ en la que hay una
    probabilidad igual de $1/2$ de que adopte el valor del número de Euler $e
    \approx 2.71828$ o el valor del número $\pi \approx 3.14159$. La entropía
    de esta distribución es $H = 1$, lo que sugiere que se pueden
    codificar los dos números mediante un dígito de un bit, $0$ o $1$,
    respectivamente. ¿Pero se puede implementar un esquema de decodificación para este código
    en una máquina de estados finitos? La respuesta es, de hecho, no, ya que se necesitan
    un número infinito de bits para describir con precisión cualquiera de los dos números.
\end{example}

Por lo tanto, por lo general es imposible disponer de un esquema de codificación y decodificación
capaz de reproducir con precisión muestras procedentes de una distribución
arbitraria de valores reales.\footnote{Es decir, si se desea codificar
    dichas muestras con precisión, la única forma es memorizar cada una de las
muestras.} Pero no tendría mucho sentido codificar una
distribución sin poder decodificarla para muestras extraídas de la
misma distribución.

Por lo tanto, para garantizar que cualquier esquema de codificación/decodificación sea computable y se pueda implementar
con memoria y recursos computacionales finitos, debemos cuantificar
la muestra $\vx$ y codificarla solo hasta una cierta precisión, digamos
$\epsilon > 0$. {\em Al hacerlo, en esencia, consideramos que dos puntos de datos cualesquiera
son equivalentes si su distancia es menor que $\epsilon$.} Más
concretamente, nos gustaría considerar esquemas de codificación
\begin{equation}
    \vx \mapsto \hat \vx
\end{equation}
de tal manera que el error esperado causado por la cuantificación esté acotado por
$\epsilon$. Es matemáticamente más conveniente, y conceptualmente
prácticamente idéntico, acotar el error esperado \textit{al cuadrado} por
\(\epsilon^{2}\), es decir,
\begin{equation}
    \Ex[d(\vx, \hat \vx)^{2}] \le \epsilon^{2}.
\end{equation}
Por lo general, se elige que la distancia \(d\) sea la distancia euclidiana, o la
norma \(L{2}\).\footnote{De manera más general, podemos sustituir \(d^{2}\) por cualquier
\textit{divergencia}.} Adoptaremos esta opción en lo sucesivo.

\subsection{La tasa-distorsión y la geometría de los datos}
Por supuesto, de entre todos los esquemas de codificación que satisfacen la restricción anterior,
nos gustaría elegir aquel que minimice la tasa de codificación resultante. Para
una variable aleatoria dada $\vx$ y una precisión $\epsilon$, esta tasa se conoce como
la \textit{tasa-distorsión}, denotada como $\cR_{\epsilon}(\vx)$.

Para comprender mejor esta tasa-distorsión, recurrimos a un teorema fundamental
de la teoría de la información, demostrado originalmente en
\textcite{shannon1959coding}, que relaciona esta tasa-distorsión con
otras magnitudes de la teoría de la información. Para comprender este teorema,
primero debemos definir una medida de \textit{diferencia} entre dos
distribuciones conocida como la divergencia de \textit{Kullback-Leibler} (KL),
o \textit{entropía relativa}, que se denota por \(\KL(\cdot \mmid
\cdot)\) y se define en relación con la entropía y la entropía cruzada
(véase la \Cref{sub:min_entropy}) como
\begin{align}
    \KL(p \mmid q) &= \CE(p \mmid q) - h(p) = \Ex_{\vx \sim p}[-\log
    q(\vx)] - \Ex_{\vx \sim p}[-\log p(\vx)] \\
    &= \Ex_{\vx \sim p}[\log (p(\vx)/q(\vx))].
\end{align}
La divergencia de KL siempre es no negativa y es igual a cero si y solo
si \(p = q\);\footnote{Técnicamente, esta igualdad debe interpretarse
    en el sentido de «casi en todas partes», es decir, salvo, posiblemente, en un conjunto de medida
    (volumen) nula, ya que este conjunto no afectaría al valor de ninguna
integral.} En el caso de que \(p\) y \(q\) tengan el mismo soporte
y ambas sean densidades de probabilidad, esto se puede demostrar mediante
el siguiente cálculo, basado en la desigualdad de Jensen y en el hecho de que
la función \(\log(\cdot)\) es estrictamente cóncava:
\begin{align}
    - \KL(p\mmid q)
    &= -\Ex_{\vx \sim p}[\log(p(\vx)/q(\vx))] = \Ex_{\vx \sim
    p}[\log(q(\vx)/p(\vx))] \\
    &\leq \log \Ex_{\vx \sim p}[q(\vx)/p(\vx)] =  \log \int_{\R^{D}}
    \frac{q(\vxi)}{p(\vxi)}p(\vxi)\odif{\vxi} \\
    &= \log \int_{\R^{D}}q(\vxi)\odif{\vxi} = \log \Ex_{\vx \sim
    q}[1] = \log 1 = 0.
\end{align}
A continuación, definimos la \textit{información mutua} \(I(\cdot; \cdot)\),
que es una medida de la \textit{dependencia} entre dos variables aleatorias, como
\begin{equation}\label{eq:mutual-information}
    I(\vx; \hat{\vx})
    = \KL(p(\vx, \hat{\vx}) \mmid p(\vx) p(\hat{\vx})),
\end{equation}
Una vez establecidas las definiciones anteriores, podemos expresar la tasa-distorsión
\(\cR_{\epsilon}\) en términos puramente probabilísticos como
\begin{equation}
    \cR_{\epsilon}(\vx)
    = \min_{p(\hat \vx \mid \vx): \Ex[\norm{\vx -\hat \vx}_2^{2}] \le
    \epsilon^{2}}
    I(\vx; \hat{\vx}).
    \label{eqn:rate-distortion-general}
\end{equation}

Tenga en cuenta que la minimización
en \eqref{eqn:rate-distortion-general} se realiza sobre todas las distribuciones condicionales
$p(\hat \vx \mid \vx)$ 
que satisfacen la restricción de distorsión $\Ex[\norm{\vx-
\hat \vx}_2^{2}] \le \epsilon^{2}$. Cada una de estas distribuciones condicionales da lugar a una distribución conjunta $p(\vx,
\hat{\vx}) = p(\hat{\vx} \mid \vx) p(\vx)$, que determina la información mutua
\eqref{eq:mutual-information}.
Muchas propiedades convenientes de la información mutua (y, por lo tanto, de la 
tasa-distorsión) se deducen de las propiedades correspondientes de la divergencia de KL.
Por la definición, sabemos que $\cR_{\epsilon}(\vx)$ es una función {\em
no creciente} en $\epsilon$.

\begin{example}[Tasa-distorsión de origen gaussiano (Teorema
    13.3.3, \cite{Cover-Thomas})]\label{ex:gaussian_rate_distortion}
    Consideremos una variable aleatoria gaussiana \(\vx \sim \dNorm(0,
    \sigma^{2})\). En este caso,
    es bien sabido que la función de tasa-distorsión se puede
    calcular en forma cerrada como:
    \begin{equation}
        \cR_{\epsilon}(x) =
        \casework{\frac{1}{2}\log\left(
            {\sigma^{2}}/{\epsilon^{2}}\right), & 0 \leq
        \epsilon \leq \sigma, \\ 0, & \epsilon > \sigma}.
    \end{equation}
    Esto se extiende de una manera particularmente interesante al caso de
    \(D\) dimensiones. Supongamos que \(\vx \sim \dNorm(\vzero,
    \vSigma)\), donde  \(\vSigma \in \PSD(D)\) tiene valores propios
    \(\lambda_{1} \geq \cdots \geq \lambda_{D} \geq 0\). Entonces, la 
    tasa-distorsión se puede obtener mediante el «relleno de agua» como
    \begin{equation}
        \cR_{\epsilon}(\vx) = \frac{1}{2} \sum_{i =
        1}^{D}\log\rp{\frac{\lambda_{i}}{\min\{\kappa, \lambda_{i}\}}},
    \end{equation}
    donde el «nivel de agua» \(\kappa\) cumple que \(\sum_{i =
    1}^{D}\min\{\kappa, \lambda_{i}\} = \epsilon^{2}\). Intuitivamente,
    las direcciones con una varianza \(\lambda_{i}\) mayor que \(\kappa\)
    se codifican de manera efectiva, mientras que las direcciones con una varianza menor
    se descartan, lo que garantiza que la restricción de distorsión se
    cumpla con una información mutua mínima. Consulte el
    \Cref{exer:KKT} para ver una derivación detallada de esta fórmula cerrada.
\end{example}

\begin{remark}
    Resulta que la tasa-distorsión es una aproximación viable
    a la entropía de $\vx$ en el siguiente sentido.
    Supongamos que $\vx$ y $\hat{\vx}$ son vectores aleatorios continuos.
    Entonces, la información mutua se puede escribir como
    \begin{equation}\label{eq:information-entropy}
        I(\vx; \hat\vx) = h(\vx) - h(\vx \mid \hat \vx),
    \end{equation}
    donde $h(\vx \mid \hat \vx) = -\bE[\log_2 p(\vx \mid \hat\vx)]$ es la
    \textit{entropía condicional} de $\vx$ dado $\hat \vx$.
    Por lo tanto, la tasa de codificación mínima se alcanza cuando la diferencia entre la
    entropía de $\vx$ y la entropía condicional de $\vx$ dada $\hat \vx$
    se minimiza entre todas las distribuciones que satisfacen
    la restricción: $\mathbb{E} [\norm{\vx- \hat \vx}_2^2] \le \epsilon^2$.

    De hecho, no es necesario suponer que $\vx$ y $\hat \vx$ sean continuos
    para llegar a una conclusión como la anterior. Por ejemplo, si ambos vectores aleatorios
    son discretos, tras una interpretación adecuada de la divergencia de KL
    para vectores aleatorios de valores discretos, tenemos que
    \begin{equation}\label{eq:mutual}
        I(\vx; \hat\vx) = H(\vx) - H(\vx \mid \hat \vx),
    \end{equation}
    donde $H(\vx)$ se define en la \eqref{eqn:entropy-discrete} de la
    \Cref{sub:min_entropy}.
    En términos más generales, se pueden emplear conceptos matemáticos avanzados de la teoría de la medida
    para definir la información mutua (y, por lo tanto, la tasa-distorsión) para
    cualquier variable aleatoria $\vx$ y $\hat \vx$, incluidas aquellas con distribuciones de baja dimensión bastante
    exóticas; véase \textcite[\S 8.5]{Cover-Thomas}.
\end{remark}

\begin{remark}
    Dado un conjunto de puntos de datos en $\vX = [\vx_1,\ldots, \vx_N]$, siempre se pueden
    interpretar como muestras de una distribución discreta uniforme con una probabilidad igual
    a $1/N$ sobre estos $N$ vectores. La entropía de dicha distribución
    es $H(\vX) = \frac{1}{N}\log_2 N$.\footnote{Nótese, una vez más, que aunque
        podamos codificar
        estos vectores con esta tasa de codificación, no podemos decodificarlos
        con una precisión
    arbitraria.} Sin embargo, aunque $\vX$ sea una distribución uniforme sobre
    sus muestras,
    la tasa de codificación $\cR_{\epsilon}(\vX)$ que se puede alcanzar con un esquema de codificación con pérdida
    podría ser significativamente inferior a $H(\vX)$ si estas muestras
    no están tan uniformemente
    distribuidas y muchas se agrupan muy cerca unas de otras. Por lo tanto, para la
    segunda distribución que se muestra en la \Cref{fig:1d-line}, si se elige adecuadamente
    el error de cuantificación  $\epsilon$, la tasa de codificación con pérdida que se puede alcanzar puede ser
    significativamente menor que si se codificara como una distribución 
    uniforme.\footnote{No obstante, para esta distribución uniforme discreta,
        cuando $\epsilon$ es lo suficientemente pequeño, siempre se cumple que $H(\vX) \approx
    \cR_{\epsilon}(\vX)$.} Obsérvese también que, con el concepto de 
    tasa-distorsión, la dificultad mencionada en el
    \Cref{eg:two-inrational} también desaparece: podemos elegir dos
    números racionales que estén lo suficientemente cerca de cada uno de los dos
    números irracionales. El esquema de codificación resultante tendrá una
    complejidad finita.
\end{remark}

\begin{example}\label{example:sphere-covering-rate-distortion}
    A veces, puede darse una situación opuesta, en la que queremos fijar primero la tasa de codificación
    e intentar encontrar un esquema de codificación que minimice la distorsión.
    Por ejemplo, supongamos que solo queremos utilizar un número fijo de códigos para
    los puntos muestreados de una distribución, y queremos saber cómo diseñar los
    códigos de tal manera que se minimice la distorsión promedio o máxima durante el
    esquema de codificación/decodificación. Por ejemplo, dada una distribución uniforme en
    un cuadrado unitario, nos preguntamos con qué precisión podemos codificar los puntos extraídos de esta
    distribución, digamos que con $n$ bits. Este problema equivale a preguntarse cuál
    es el radio mínimo (es decir, la distorsión) tal que podamos cubrir el cuadrado
    unitario con $2^n$ discos de ese radio.
    La \Cref{fig:seven-circles-packing} muestra cubrimientos aproximadamente óptimos de
    un cuadrado con \(n = 4, 6, 8\), de modo que \(2^{n} = 16, 64, 256\)
    discos, respectivamente. Obsérvese que los radios óptimos de los discos
    disminuyen a medida que aumenta el número de discos \(2^{n}\).
    \begin{figure}
        \centering
        \includegraphics[width=0.3\textwidth]{\toplevelprefix/chapters/chapter4/figs/16.png}
        \hfill
        \includegraphics[width=0.3\textwidth]{\toplevelprefix/chapters/chapter4/figs/64.png}
        \hfill
        \includegraphics[width=0.3\textwidth]{\toplevelprefix/chapters/chapter4/figs/256.png}

        \caption{Aproximaciones a las soluciones óptimas para \(2^{4}\),
            \(2^{6}\) y \(2^{8}\) discos que cubren un cuadrado, junto con los
            radios correspondientes, calculados mediante un algoritmo
        de optimización heurística.}
        \label{fig:seven-circles-packing}
    \end{figure}
\end{example}

\begin{figure}[t]
    \centering
    \includegraphics[width=0.6\textwidth]{\toplevelprefix/chapters/chapter4/figs/continuation.png}
    \caption{\small\textbf{La aproximación de una distribución de baja dimensión
        mediante \(\epsilon\) bolas.} Podemos observar que, a medida que el parámetro \(\epsilon\)
        se reduce, la unión de las bolas de radio \(\epsilon\) se aproxima cada vez mejor al soporte de la
        distribución real (negra). Además, los
        filtros de eliminación de ruido asociados (cuya correspondencia entre entrada y salida viene dada por las flechas proporcionadas)
        obtenidos al aproximar la distribución real mediante una mezcla
        de distribuciones gaussianas,
        cada una con covarianza \((\epsilon^{2}/D)\vI\), se acercan cada vez
        más
        a los filtros de eliminación de ruido verdaderos.
        El \Cref{thm:covering-number-rate-distortion}
        establece que esta aproximación caracteriza la función de tasa-distorsión
        para distorsiones pequeñas $\epsilon$, unificando los
        enfoques paralelos
        de minimización de la tasa de codificación y eliminación de ruido para el aprendizaje de distribuciones de baja dimensión
    sin patologías.}
    \label{fig:continuation}
\end{figure}

Resulta ser un problema notoriamente difícil obtener expresiones en forma cerrada
para la función de tasa-distorsión \eqref{eqn:rate-distortion-general} para distribuciones generales
$p(\vx)$. Sin embargo, tal y como
sugiere el \Cref{example:sphere-covering-rate-distortion}, existen casos especiales importantes
en los que la \textit{geometría} del soporte de la distribución
$p(\vx)$ puede relacionarse con la función de tasa-distorsión y, por lo tanto, con la tasa de codificación óptima
para un nivel de distorsión de $\epsilon$. De hecho, este ejemplo puede generalizarse a cualquier contexto en el que el
soporte de $p(\vx)$ sea un conjunto compacto suficientemente regular —incluidas
las distribuciones de baja dimensión— y $p(\vx)$ tenga una distribución uniforme en su
soporte.
Esto abarca una gran cantidad de casos de interés práctico.
Formalizamos esta noción en el siguiente resultado, que establece esta
propiedad para un caso especial.

\begin{theorem}\label{thm:covering-number-rate-distortion}
    Supongamos que \(\vx\) es una variable aleatoria tal que su soporte
    \(K \doteq \Supp(\vx)\) es un conjunto compacto. Definamos el
    \textit{número de cobertura} \(\cN_{\epsilon}(K)\) como el número mínimo
    de bolas de radio \(\epsilon\) que pueden cubrir \(K\), es decir,
    \begin{equation}
        \cN_{\epsilon}(K) \doteq \min\left\{n \in \bN \colon \exists
            \vp_{1}, \dots, \vp_{n} \in K\ \text{s.t.}\ K \subseteq
        \bigcup_{i = 1}^{n}B_{\epsilon}(\vp_{i})\right\},
    \end{equation}
    donde \(B_{\epsilon}(\vp) = \{vxi \in \R^D \mid \norm{\vxi -
            \vp}_2 \leq
    \epsilon\}\) es la bola euclidiana de radio \(\epsilon\) centrada
    en \(\vp\).
    Entonces se cumple que
    \begin{equation}
        \cR_{\epsilon}(\vx)
        \leq \log_{2} \cN_{\epsilon}(K).
    \end{equation}
    Si, además, $\vx$ sigue una distribución uniforme en $K$ y $K$ es
    una unión de hiperesferas de baja dimensión mutuamente ortogonales,\footnote{De hecho, es
    posible tratar casos de $K$ altamente irregulares, como variedades compactas o incluso fractales,
    con un resultado paralelo,
    pero su formulación
    se vuelve mucho más técnica: véase Riegler et al.
    \cite{Riegler2018-jh,Riegler2023-rr}. En el
    \Cref{app:rate-distortion-covering} ofrecemos una demostración sencilla que muestra el resultado
    para uniones de hiperesferas ortogonales, una suposición que equivale a
    una mezcla de distribuciones gaussianas de rango bajo con soportes ortogonales cuando la dimensión
    de cada subespacio es suficientemente grande. Véase el
    \Cref{app:rate-distortion-covering} para un análisis más detallado de esta
    hipótesis.}
    entonces se cumple una cota inferior correspondiente:
    \begin{equation}
        \cR_{\epsilon}(\vx)
        \geq
        \log_{2} \cN_{\epsilon}(K) - O(D).
    \end{equation}
\end{theorem}
\begin{proof}
    La demostración de este teorema excede el alcance de este libro, por lo que
    la remitimos al \Cref{app:rate-distortion-covering}.
\end{proof}

La cota superior del \Cref{thm:covering-number-rate-distortion} solo requiere que
$\Supp(\vx)$ sea un conjunto compacto, mientras que nuestra demostración de la cota inferior requiere
supuestos de regularidad más fuertes: concretamente, que $\vx$ esté uniformemente
distribuido en $\Supp(\vx)$ y que $\Supp(\vx)$ sea una mezcla de
gaussianas de rango bajo (con la normalización adecuada).\footnote{Más adelante en este capítulo, 
en la \Cref{subsec:MCR2}, veremos una función objetivo cuyos máximos son mezclas normalizadas 
de gaussianas de bajo rango que se corresponden con los datos de entrada de cierta manera.
El \Cref{ch:deep-networks} desarrollará aún más este marco para el aprendizaje profundo,
demostrando que la suposición estructural que hacemos para demostrar que la cota inferior
no es excesivamente restrictiva.}
Bajo estos supuestos, la implicación de el
\Cref{thm:covering-number-rate-distortion} puede resumirse de la siguiente manera: para
una codificación suficientemente precisa de la distribución de
$\vx$ (es decir, cuando $\epsilon \to 0$), el marco de codificación de mínima tasa-distorsión
es equivalente al
problema del empaquetamiento de bolas en el soporte de $\vx$.
Así pues, la tasa-distorsión puede considerarse como una forma «sensible a la probabilidad» de
aproximar el soporte de la distribución de 
\(\vx\) mediante una mezcla de muchas bolas pequeñas (\Cref{fig:continuation}).

Ahora analizaremos otra relación entre esto y la
jerarquía de complejidad «ruido-difusión-entropía» que hemos comentado anteriormente en este
capítulo.

\begin{remark}\label{rem:slb}
    El ingrediente clave en la demostración de la cota inferior en el
    \Cref{thm:covering-number-rate-distortion} es un resultado importante de
    la teoría de la información conocido como el \textit{límite inferior de Shannon} para la 
    tasa-distorsión, que lleva el nombre de Claude Shannon, quien lo derivó por primera vez en un
    caso especial
    \cite{shannon1959coding}. En él
    se establece la siguiente estimación para la función de tasa-distorsión,
    para cualquier variable aleatoria
    $\vx$ con
    una densidad $p(\vx)$ y una norma al cuadrado que se espera finita \cite{Linder1994-ej}:
    \begin{equation}\label{eq:slb}
        \cR_{\epsilon}(\vx)
        \geq
        h(\vx)
        - \log_2 \volume(B_{\epsilon})
        -
        C_D,
    \end{equation}
    donde $C_D > 0$ es una constante que depende únicamente de $D$. Además,
    esta cota inferior
    es, de hecho, exacto cuando $\epsilon \to 0$: es decir,
    \begin{equation}
        \lim_{\epsilon \to 0} \cR_{\epsilon}(\vx) - \left[ h(\vx) - \log_2
        \volume(B_{\epsilon}) - C_D\right] = 0.
    \end{equation}
    Así pues, cuando la distorsión $\epsilon$ es pequeña, podemos pensar únicamente en términos de la
    cota inferior de Shannon, en lugar del problema de optimización (generalmente inmanejable)
    que define la relación tasa-distorsión
    \eqref{eqn:rate-distortion-general}.

    La cota inferior de Shannon constituye el nexo entre la tasa de codificación, la minimización de la entropía,
    la eliminación de ruido y los enfoques de empaquetamiento geométrico de bolas para el aprendizaje
    de distribuciones de baja dimensión. Obsérvese que, en el caso especial de una densidad uniforme
    $p(\vx)$, \eqref{eq:slb} se convierte en
    \begin{align}
        \cR_{\epsilon}(\vx) &\geq -\int_K \frac{1}{\volume(K)} \log_2
        \frac{1}{\volume(K)} \odif \vxi
        - \log_2 \volume(B_{\epsilon}) - C_D
        \\
        &=
        \log_2 \volume(K) / \volume(B_{\epsilon})
        - C_D. \label{eq:shannon-lower-bound-uniform}
    \end{align}
    La relación $\volume(K) / \volume(B_{\epsilon})$ aproxima el número de
    bolas de radio $\epsilon$ necesarias para cubrir $K$ mediante un argumento del peor caso, lo cual es
    preciso para conjuntos $K$ suficientemente regulares cuando $\epsilon$ es pequeño (véase
    el \Cref{app:rate-distortion-covering} para más detalles).
    Mientras tanto, recordemos el modelo gaussiano de eliminación de ruido $\vx_{\epsilon} = \vx
    + \epsilon \vg$ mencionado
    anteriormente en este capítulo, donde $\vg \sim \cN(\vzero, \vI)$ es independiente de
    $\vx$.
    Curiosamente, la entropía diferencial de la distribución conjunta $(\vx,
    \epsilon\vg)$
    se puede calcular como
    \begin{align}
        h(\vx, \epsilon\vg)
        &=
        -\int
        p(\vxi) p(\vgamma) \log_2 p(\vxi) p(\vgamma) \odif \vxi \odif \vgamma
        \\
        &=
        h(\vx) + h(\epsilon \vg).
    \end{align}
    Hemos visto que la entropía gaussiana calculada en la
    \Cref{eqn:entropy-Gaussian-multi}: cuando $\epsilon$ es pequeño, es igual, salvo
    por constantes aditivas, a la cantidad $-\log_2
    \volume(B_{\epsilon})$ que hemos visto en la cota inferior de Shannon.
    En ciertos casos especiales (por ejemplo, datos apoyados en subespacios incoherentes
    de bajo rango),
    cuando $\epsilon$ es pequeño y el soporte de $\vx$ es suficientemente regular,
    la distribución de $\vx_{\epsilon}$ puede incluso aproximarse bien
    localmente por el
    producto de las distribuciones $p(\vx)$ y $p(\vg)$, lo que justifica el cálculo anterior.
    Por lo tanto, el proceso de eliminación de ruido gaussiano ofrece otra interpretación más de la
    cota inferior de Shannon, ya que se deriva de la entropía de una versión con ruido de
    $\vx$, con un nivel de ruido proporcional al nivel de distorsión $\epsilon$.

    Por lo tanto, este enfoque de distorsión de tasa finita mediante el recubrimiento de bolas reactiva o
    generaliza todas las medidas anteriores de complejidad de la
    distribución, lo que nos permite
    diferenciar y clasificar las distintas distribuciones de manera unificada.
    Estos puntos de vista interrelacionados se visualizan en la \Cref{fig:continuation}.
\end{remark}

En el caso de una distribución general con niveles de distorsión finitos, suele
ser imposible hallar su función de tasa-distorsión
en forma analítica. A menudo hay que recurrir al
cálculo numérico\footnote{Los lectores interesados pueden consultar
    \cite{Blahut-1972} para conocer un algoritmo clásico que calcula numéricamente la
función de tasa-distorsión para una distribución discreta.}.
Sin embargo, como veremos, en nuestro contexto a menudo necesitamos conocer
la tasa-distorsión como una función explícita de un conjunto de puntos de datos
o de sus representaciones. Esto se debe a que queremos utilizar la tasa de codificación
como medida de la calidad de las representaciones. Por lo tanto, deberíamos trabajar con
distribuciones cuyas funciones de tasa-distorsión adopten formas analíticas
explícitas. Con este fin, comenzaremos con la familia de distribuciones más simple y, al mismo tiempo,
la más importante.

\subsection{Tasa de codificación con pérdida para una distribución gaussiana
de baja dimensión}\label{subsec:lossy DR}
Supongamos ahora que se nos da un conjunto de muestras de datos en $\vX = [\vx_1,
\ldots, \vx_N]$ procedente de cualquier distribución.\footnote{O bien, estos puntos de datos
podrían considerarse ellos mismos como una distribución (empírica).} Nos gustaría
elaborar un esquema constructivo que pueda codificar los datos
con cierta precisión, digamos
\begin{equation}
    \vx_i \mapsto \hat \vx_i, \quad \mbox{subject to} \quad \|\vx_i -
    \hat \vx_i\|_2 \le \epsilon.
\end{equation}
Nótese que se trata de una condición suficiente, explícita e interpretable que
garantiza que los datos se codifican de tal manera que \(\frac{1}{N}\sum_{i = 1}^{N}
\norm{\vx_i- \hat \vx_i}_2^{2} \le \epsilon^{2}\). Esta última
desigualdad es precisamente la restricción de tasa-distorsión para la distribución
empírica dada y su codificación. Por ejemplo, en el
\Cref{example:sphere-covering-rate-distortion}, utilizamos este
criterio simplificado para determinar explícitamente la distorsión mínima y
el esquema de codificación explícito para una tasa de codificación dada.

Sin pérdida de generalidad, supongamos que la media de $\vX$ es cero,
es decir, $\frac{1}{N} \sum_{i = 1}^{N} \vx_i = \boldsymbol{0}$. Sin
ningún conocimiento previo sobre la naturaleza de la distribución subyacente a $\vX$,
podemos considerar que $\vX$ se ha obtenido a partir de una distribución gaussiana
$\mathcal{N}(\boldsymbol{0}, {\boldsymbol{\Sigma}})$ con la
covarianza\footnote{Se sabe que, dada una varianza fija, la
    distribución gaussiana alcanza la máxima entropía. Es decir, establece un límite superior
para el peor de los casos en cuanto a la tasa de codificación posible.}
\begin{equation}
    {\boldsymbol{\Sigma}} = \frac{1}{N} \vX\vX^\top.
\end{equation}
Obsérvese que, desde el punto de vista geométrico, ${\boldsymbol{\Sigma}}$ caracteriza una
región elipsoidal en la que se encuentran la mayoría de las muestras $\vx_i$.

Podemos considerar $\hat \vX = [\hat \vx_1,\ldots, \hat \vx_N]$ como una versión con ruido
de $\vX = [\vx_1, \ldots, \vx_N]$:
\begin{equation}
    \hat \vx_i = \vx_i + \vw_i,
\end{equation}
donde $\vw_i  \sim
\mathcal{N}(\boldsymbol{0} , {\epsilon^2}  \boldsymbol{I}/{D})$ es un
ruido gaussiano
independiente de $\vx_i$. Entonces, la covarianza de $\hat \vx_i$ viene dada por
\begin{equation}
    \hat{\boldsymbol{\Sigma}} = \mathbb{E}\left[\hat \vx_i \hat
    \vx_i^\top\right] = \frac{\epsilon^2}{D} \boldsymbol{I} +
    \frac{1}{N} \vX\vX^\top.
\end{equation}
Tenga en cuenta que el volumen de la región delimitada por los vectores $\vx_i$ es
proporcional a la raíz cuadrada del determinante de la
matriz de covarianza
\begin{equation}
    \mbox{volume}(\hat \vx_i) \propto \sqrt{\det
    \big(\hat{\boldsymbol{\Sigma}}\big)} =
    \sqrt{\det\left(\frac{\epsilon^2}{D} \boldsymbol{I} + \frac{1}{N}
    \vX\vX^\top \right)}.
\end{equation}
El volumen abarcado por cada vector aleatorio $\boldsymbol{w}_i$ es
proporcional a
\begin{equation}
    \mbox{volume}(\boldsymbol{w}_i) \propto
    \sqrt{\det\left(\frac{\epsilon^2}{D} \boldsymbol{I} \right)}.
\end{equation}

\begin{figure}
    \centering
    \includegraphics[width=\linewidth]{\toplevelprefix/chapters/chapter4/figs/Gaussian-compression.png}
    \caption{Cubrir la región abarcada por los vectores de datos utilizando
        bolas de radio $\epsilon$. Cuanto mayor sea el volumen del espacio, más
        bolas se necesitarán y, por lo tanto, más bits se necesitarán para codificar y
        enumerar las bolas. Cada vector de valor real $\vx$ puede codificarse
    como el número de la bola en la que se encuentra.}
    \label{fig:ball-packing}
\end{figure}

Para codificar los vectores que se encuentran dentro de la región definida por $\hat \vx_i$,
podemos cubrir dicha región con bolas no superpuestas
de radio $\epsilon$, tal como se ilustra en la
\Cref{fig:ball-packing}. Cuando el volumen de la región abarcada por
$\hat \vx_i$ es significativamente mayor que el volumen de la
bola de radio $\epsilon$, el número total de bolas que necesitamos para cubrir la región es
aproximadamente igual a la relación entre ambos volúmenes:
\begin{equation}
    \# \,\epsilon\mbox{-balls} \approx \frac{\mbox{volume}(\hat
    \vx_i)}{\mbox{volume}(\vw_i)} = \sqrt{\det\left(\boldsymbol{I} +
    \frac{D}{N\epsilon^2} \vX\vX^\top  \right)}.
\end{equation}
Si utilizamos números binarios para etiquetar todas las bolas de radio $\epsilon$ en la región de
interés, el número total de bits binarios necesarios es, por lo tanto, %
\begin{equation}
    \cR_{\epsilon}(\vX) \approx %
    \log_2 (\# \,\epsilon\mbox{-balls}) \approx R_{\epsilon}(\vX)
    \doteq \frac{1}{2} \log \det \left(\boldsymbol{I} +
    \frac{D}{N\epsilon^2} \vX\vX^\top \right).
    \label{eqn:rate-Gaussian}
\end{equation}
Nótese que, a partir del \Cref{ex:gaussian_rate_distortion}, se puede mostrar que
el lado derecho de esta fórmula es igual a la 
tasa-distorsión (con una precisión de \(\epsilon\)) de una fuente gaussiana con
matriz de covarianza \(\frac{1}{N} \vX\vX^\top + (\epsilon^{2}/D) \vI\)
(por ejemplo, la matriz de covarianza empírica de los datos más un
término de ruido). De este modo, un argumento geométrico sobre el volumen de los
vectores encerrados y el empaquetamiento de esferas se traduce en un
argumento basado en la teoría de la información sobre la tasa-distorsión de una
fuente gaussiana perturbada.

\begin{example}
    La \Cref{fig:ball-packing} muestra un ejemplo de una distribución bidimensional
    con un soporte elipsoidal, que se aproxima al soporte de una distribución gaussiana
    bidimensional. La región está cubierta por pequeñas bolas de
    tamaño $\epsilon$. Todas las bolas están numeradas del $1$ al, digamos, $n$.
    Entonces, dado cualquier vector $\vx$ en esta región,
    solo necesitamos determinar a qué centro de bolas de $\epsilon$ está
    más cerca, lo que se denota como $\operatorname{ball}_{\epsilon}(\vx)$. Para
    recordar $\vx$, solo necesitamos recordar el número de esta bola,
    cuyo almacenamiento requiere $\log(n)$ bits. Si necesitamos decodificar $\vx$
    a partir de este número, simplemente tomamos $\hat \vx$ como el centro de la
    bola. Esto da lugar a un esquema explícito de codificación y decodificación:
    \begin{equation}
        \vx \longrightarrow \operatorname{ball}_{\epsilon}(\vx)
        \longrightarrow \hat \vx = \mbox{center of}
        \operatorname{ball}_{\epsilon}(\vx).
    \end{equation}
    Se puede decir que estos centros de las bolas son los «códigos» de un libro de códigos
    o de un diccionario para el esquema de codificación. Es fácil ver que
    la precisión de este esquema de codificación-decodificación (con pérdida) es aproximadamente
    el radio de la bolas $\epsilon$. 
    Es evidente que $\cR_{\epsilon}(\vZ)$ es el número promedio
    de bits necesarios para codificar el número de bola de cada vector $\vz$ con
    este esquema de codificación y, por lo tanto, puede denominarse la {\em tasa de codificación}
    asociada a este esquema.
\end{example}

A partir de la derivación anterior, sabemos que la tasa de codificación
$\cR_{\epsilon}(\vX)$ es (aproximadamente) alcanzable mediante un esquema explícito
de codificación (y decodificación). Presenta dos propiedades interesantes:
\begin{itemize}
    \item En primer lugar, se puede observar que $R_{\epsilon}(\vX)$ se asemeja mucho
        a la función de tasa-distorsión de una fuente gaussiana
        \cite{Cover-Thomas} (véase el \Cref{ex:gaussian_rate_distortion}). De hecho,
        cuando $\epsilon$ es pequeño, la expresión anterior es una aproximación
        muy cercana a la tasa-distorsión de una fuente gaussiana, tal como
        se señala en   \cite{MaY2007-PAMI}. 
    \item En segundo lugar, la misma tasa de codificación en forma cerrada $R_{\epsilon}(\vX)$
        puede derivarse como una aproximación de \(\cR_{\epsilon}(\vX)\)
        si se supone que los datos $\vX$ proceden de un subespacio lineal.
        Esto se puede demostrar cuantificando adecuadamente la descomposición en valores singulares
        (SVD) de $\vX = \boldsymbol{U}
        \boldsymbol{\Sigma}\boldsymbol{V}^\top$ y construyendo un
        esquema de codificación con pérdidas para vectores en el subespacio generado por
        $\boldsymbol{U}$ \cite{MaY2007-PAMI}. 
\end{itemize}
En nuestro contexto, la expresión en forma cerrada $R_{\epsilon}(\vX)$ es
bastante fundamental: se trata de la tasa de codificación asociada a un esquema de codificación con pérdidas explícito
y natural para datos extraídos de una distribución gaussiana
o de un subespacio lineal. Como veremos en el próximo
capítulo, esta fórmula desempeña un papel importante en la comprensión de la
arquitectura de las redes neuronales profundas.

\subsection{Agrupamiento de una mezcla de gaussianas de baja dimensión}
\label{sec:clustering-Gaussians}
Como ya hemos comentado anteriormente, el conjunto de datos dado $\vX$ suele tener
estructuras intrínsecas de baja dimensión. Por lo tanto, codificarlo como una distribución
gaussiana general resultaría muy redundante. Si logramos identificar esas estructuras intrínsecas
en $\vX$, podríamos diseñar esquemas de codificación mucho mejores que
ofrezcan tasas de codificación mucho más bajas. O, dicho de otra forma, los códigos utilizados para
codificar dicho $\vX$ pueden comprimirse. Veremos que la compresión
ofrece una forma computable y unificadora de identificar tales estructuras. En esta
sección, ilustramos esta importante idea con la familia más básica
de estructuras de baja dimensión: una mezcla de gaussianas (de baja dimensión)
o subespacios.

\begin{example}
    La \Cref{fig:two-subspaces} muestra un ejemplo en el que los datos $\vX$
    se distribuyen en torno a dos subespacios (o gaussianas de baja dimensión).
    Si se consideran y se codifican conjuntamente como una sola
    gaussiana, el libro de códigos discreto (con pérdida) asociado, representado
    por todas las bolas azules, es obviamente muy redundante. Podemos intentar
    identificar las ubicaciones de los dos subespacios, denotados por $S_1$ y
    $S_2$ y diseñar un libro de códigos que abarque únicamente esos dos subespacios,
    es decir, las bolas verdes. Si logramos dividir correctamente las muestras de
    los datos $\vX$ en dos subespacios: $\vX = [\vX_1, \vX_2]\bm \Pi$
    con $\vX_1 \in S_1$ y $\vX_2 \in S_2$, donde $\bm \Pi$ denota una
    matriz de permutación, entonces la tasa de codificación resultante para los datos
    será mucho menor. Esto proporciona una representación más parsimoniosa y, por lo tanto, más
    deseable de los datos.
\end{example}

\begin{figure}
    \centering
    \includegraphics[width=0.5\linewidth]{\toplevelprefix/chapters/chapter4/figs/Two-subspaces.png}
    \caption{Comparación de dos esquemas de codificación con pérdidas para datos que se
        distribuyen alrededor de dos subespacios. Uno consiste en agrupar (azul)
        bolas de radio $\epsilon$ para todo el espacio generado por los dos
        subespacios; el otro consiste en agrupar bolas solo en una vecindad
        tabular alrededor de los dos subespacios. Este último, obviamente, tiene un
        libro de códigos mucho más pequeño y da como resultado una tasa de codificación mucho menor
    para las muestras en los subespacios.}
    \label{fig:two-subspaces}
\end{figure}

Por lo tanto, en términos más generales, si los datos proceden de cualquier mezcla
de subespacios o gaussianas de baja dimensión, sería conveniente
identificar esos componentes y codificar los datos basándose en las dimensiones intrínsecas
de dichos componentes. Resulta que no perdemos mucha
generalidad al suponer que los datos provienen de una mezcla de
distribuciones gaussianas de baja dimensión. Esto se debe a que una mezcla de distribuciones gaussianas puede
aproximar con gran precisión a la mayoría de las distribuciones generales \cite{borkar2016gaussian}.

\paragraph{El problema de agrupamiento.}
Ahora bien, en el caso de esta familia concreta de distribuciones, ¿cómo podemos identificar de manera eficaz
y eficiente esos componentes de baja dimensión a partir de un conjunto de muestras
\begin{equation}
    \vX = \left[\vx_1, \vx_2, \ldots, \vx_N\right],
\end{equation}
extraído de ellos? En otras palabras, dado el conjunto de datos completo $\vX$,
queremos dividirlo, o agruparlo, en varios subconjuntos, digamos $K$:
\begin{equation}
    \vX\bm \Pi = [\vX_1, \vX_2, \dots, \vX_K],
\end{equation}
donde cada subconjunto está formado por muestras extraídas de una sola
gaussiana o subespacio de baja dimensión y $\bm \Pi$ es una matriz de permutación
que indica la pertenencia a la partición. Cabe señalar que, dependiendo de
la situación, la partición podría ser determinista o
probabilística. Como se muestra en \cite{ma2007segmentation}, en el caso de una mezcla de
distribuciones gaussianas, la partición probabilística no da lugar a una menor tasa de codificación.
Por lo tanto, en aras de la simplicidad, aquí consideramos únicamente una partición determinista.

\paragraph{Agrupamiento mediante compresión con pérdidas.}
La principal dificultad para resolver el problema de agrupamiento mencionado anteriormente radica en que
normalmente no conocemos el número de agrupamientos $K$, ni tampoco
la dimensión de cada componente. El estudio de este problema de agrupamiento
cuenta con una larga trayectoria. El libro de texto \cite{GPCA} ofrece
una exposición sistemática y exhaustiva de los diferentes enfoques de
este problema. Para encontrar un enfoque eficaz para este problema, primero
debemos comprender y aclarar por qué queremos agrupar los datos. En otras
palabras, ¿qué ganamos exactamente al agrupar los datos, en comparación
con no hacerlo? ¿Cómo medimos esa ganancia? Desde el punto de vista de la
compresión de datos, una agrupación correcta debería dar lugar a un esquema de
codificación (y decodificación) más eficiente.

Para cualquier conjunto de datos dado $\vX$, ya existen dos esquemas de codificación
obvios que sirven de referencia. Estos representan dos formas extremas de codificar los datos:
\begin{itemize}
    \item Basta con considerar todas las muestras juntas como si procedieran de una
        única distribución gaussiana. La tasa de codificación asociada, tal y como se ha
        deducido anteriormente, viene dada por: 
        \begin{equation}
            \cR_{\epsilon}(\vX) \approx R_{\epsilon}(\vX) =
            \frac{1}{2} \log \det \left(\boldsymbol{I} +
            \frac{D}{N\epsilon^2} \vX\vX^\top \right).
        \end{equation}
    \item Basta con memorizar todas las muestras por separado asignando un
        número diferente a cada una. La tasa de codificación sería: 
        \begin{equation}
            \cR_0(\vX) = \log(N).
        \end{equation}
\end{itemize}

Tenga en cuenta que cualquiera de los dos esquemas de codificación puede convertirse en la solución «óptima»
para una determinada elección (extrema) del error de cuantificación $\epsilon$:
\begin{enumerate}
    \item {\em Régimen perezoso}: Si elegimos un valor de $\epsilon$ extremadamente
        grande, todas las muestras de $\vX$ pueden quedar cubiertas por una sola bola.
        La tasa es  $\lim_{\epsilon \rightarrow \infty} \cR_\epsilon
        \rightarrow \frac{1}{2}\log\det (\boldsymbol{I}) = 0$.
    \item {\em Régimen de memorización}: Si $\epsilon$ es extremadamente
        pequeño, cada muestra de $\vX$  queda cubierta por una
        bola de $\epsilon$ diferente, por lo que el total es $N$. La tasa es
        $\lim_{\epsilon \rightarrow 0} \cR_\epsilon \rightarrow \log(N)$. 
\end{enumerate}
Tenga en cuenta que el primer esquema corresponde al caso en el que
no nos interesa en absoluto ningún aspecto relevante de la distribución.
No queremos reservar ningún bit para información relevante. A esto lo llamamos
el «régimen perezoso». El segundo esquema corresponde al
caso en el que se desea decodificar cada muestra con una precisión extremadamente alta.
Por lo tanto, sería mejor «memorizar» cada muestra. A esto lo llamamos
el «régimen de memorización».

Nótese que estas asintóticas en términos de \(\epsilon\) son exactamente
las mismas que se analizaron al presentar la tasa de codificación con pérdida como una magnitud
que los modelos de difusión entrenados de forma realista aproximan, tal como se muestra al
final de la \Cref{sec:diffusion_memorization}. En la situación actual,
codificamos \textit{explícitamente} la distribución con una precisión de \(\epsilon\); en
el caso de los modelos de difusión, codificamos \textit{implícitamente} la
distribución. Ciertas propiedades del esquema, como la memorización
y la generalización, se manifiestan de manera similar en ambos casos y, de hecho,
se puede estudiar un caso para comprender el otro.

\begin{figure}
    \centering
    \includegraphics[width=0.9\linewidth]{\toplevelprefix/chapters/chapter4/figs/circle-packing.png}
    \caption{Una serie de muestras aleatorias en un plano bidimensional. Consideremos un
        disco de radio $\epsilon$ asociado a cada muestra, con la muestra como
    centro. La densidad de las muestras aumenta de izquierda a derecha.}
    \label{fig:circle-packing}
\end{figure}
\begin{example}
    Para determinar cuándo es preferible o no el régimen de memorización, consideremos
    un número, digamos $N$, de muestras distribuidas aleatoriamente en un
    área unitaria sobre un plano bidimensional.\footnote{Supongamos que los puntos se generan mediante un
    proceso de Poisson con una densidad de $N$ puntos por área unitaria.} Imaginemos
    que intentamos diseñar un esquema de codificación con pérdida con un error de cuantificación
    fijo $\epsilon$. Esto equivale a colocar un
    disco de $\epsilon$ alrededor de cada muestra, como se muestra en la
    \Cref{fig:circle-packing}. Cuando $N$ es pequeño, la probabilidad de que todos
    los discos se superpongan entre sí es nula. En este caso, es necesario y óptimo un libro de códigos de tamaño $N$.
    Cuando $N$ o la densidad
    alcanza un determinado valor crítico $N_c$, con alta probabilidad todos
    los discos comienzan a superponerse y a conectarse en un único grupo que
    cubre todo el plano; este fenómeno se conoce como «percolación» en el continuo
    \cite{Gilbert-1961,Mertens-Moore-2012}. Cuando $N$
    supera este valor, los discos se superponen en gran medida. El
    número $N$ de discos resulta muy redundante, ya que solo queremos
    codificar puntos en el plano con la precisión dada
    $\epsilon$. El número de discos necesarios para cubrir todas las muestras
    es muy inferior a $N$.\footnote{De hecho, existen algoritmos eficientes
    para encontrar dicha cobertura \cite{Booth-2001}.}
\end{example}

Tanto el régimen «perezoso» como el de memorización son en cierta medida triviales y
quizás tengan poco interés teórico o práctico. Cualquiera de
estos dos esquemas estaría lejos de ser óptimo si se utilizara para codificar un gran número
de muestras extraídas de una distribución que tenga un {\em soporte compacto y
de baja dimensión}. El régimen interesante se encuentra entre estos dos.
\begin{example}
    \begin{figure}[t]
        \centering
        \includegraphics[width=0.4\linewidth]{\toplevelprefix/chapters/chapter4/figs/Two-lines-and-plane.png}
        \includegraphics[width=0.9\linewidth]{\toplevelprefix/chapters/chapter4/figs/Coding-Rate.jpg}
        \caption{Parte superior: 358 muestras ruidosas extraídas de dos rectas y un
            plano en $\R^3$. Parte inferior: el efecto de variar
            $\epsilon$ sobre el resultado de la agrupación y la tasa de codificación. La
            línea roja marca la varianza $\epsilon_0$ del ruido gaussiano
            añadido a las muestras. Obsérvese que los gráficos de la parte inferior izquierda y
            del centro muestran la misma magnitud con diferentes
        ubicaciones y escalas de los ejes.}
        \label{fig:two-lines-and-plane}
        \label{fig:two-lines-and-plane-epsilon}
    \end{figure}
    La \Cref{fig:two-lines-and-plane} muestra un ejemplo con muestras ruidosas
    extraídas de dos rectas y un plano en $\R^3$. Como se observa en el
    gráfico (c) de la derecha, la tasa de codificación óptima disminuye de forma monótona a medida que
    aumentamos $\epsilon$, tal y como se esperaba a partir de la propiedad de la función de tasa-distorsión.
    Los gráficos (a) y (b) muestran, al variar $\epsilon$ desde valores muy
    pequeños (cercanos a cero) hasta valores muy grandes (cercanos al infinito), el número óptimo de
    grupos cuando la tasa de codificación es mínima. Se pueden observar claramente el régimen «perezoso»
    y el régimen de memorización en los dos extremos de los gráficos. Pero también se puede
    observar en el gráfico (b) que, cuando el error de cuantificación $\epsilon$ se elige
    en torno al nivel de la varianza real del ruido $\epsilon_0$, el número óptimo
    de agrupaciones es el número «correcto» tres, que representa dos líneas y
    un plano.  Nos referimos informalmente a este régimen intermedio como el
    «régimen de generalización». Obsérvese que se produce una transición de fase brusca
    entre estos regímenes.\footnote{Hasta donde sabemos,
        no existe una justificación teórica rigurosa para estos
    comportamientos de transición de fase.}
\end{example}

A partir del análisis y los ejemplos anteriores, vemos que, cuando el
error de cuantificación en relación con la densidad de la muestra\footnote{o la
densidad de la muestra en relación con el error de cuantificación} se encuentra dentro de un
rango adecuado, minimizar la tasa de codificación con pérdida nos permitiría descubrir
la distribución subyacente (de baja dimensión) de los datos muestreados.
{\em Por lo tanto, la cuantificación, que comenzó como una elección de carácter práctico, parece
    estar convirtiéndose en algo necesario para aprender una distribución continua a partir de
su distribución empírica con un número finito de muestras.} Aunque aún no se ha logrado establecer una teoría rigurosa
que explique este fenómeno, aquí, con
fines de aprendizaje, nos interesa cómo aprovechar este fenómeno para
diseñar algoritmos capaces de encontrar la distribución correcta.

Utilicemos el sencillo ejemplo que se muestra en la \Cref{fig:two-subspaces} para
ilustrar las ideas básicas. Si se pueden dividir todas las muestras en $\vX$
en dos grupos en $\vX_1$ y $\vX_2$, con $N_1$ y $N_2$ muestras
respectivamente, entonces la tasa de codificación asociada sería\footnote{Aquí
    ignoramos algunos bits de sobrecarga necesarios para codificar la pertenencia de
cada muestra, por ejemplo, mediante la codificación de Huffman.}
\begin{equation}
    R_{\epsilon}^c(\vX\mid \boldsymbol{\Pi}) =
    \frac{N_1}{N}R_{\epsilon}(\vX_1) + \frac{N_2}{N}R_{\epsilon}(\vX_2),
\end{equation}
donde utilizamos $\boldsymbol{\Pi}$ para indicar la pertenencia a la
partición. Si la partición respeta las estructuras de baja dimensión
de la distribución —en este caso, la pertenencia de $\vX_1$ y $\vX_2$ a los
dos subespacios, respectivamente—, entonces la tasa de codificación resultante debería ser
significativamente menor que la de los dos esquemas básicos anteriores:
\begin{equation}
    R_{\epsilon}^c(\vX \mid \boldsymbol{\Pi}) \ll R_{\epsilon}(\vX),
    \quad     R_{\epsilon}^c(\vX \mid \boldsymbol{\Pi}) \ll R_0(\vX).
\end{equation}
En general, podemos reformular el problema de agrupamiento como un problema de optimización
que minimiza la tasa de codificación:
\begin{equation}
    \min_{\boldsymbol{\Pi}}  \bc{ R_{\epsilon}^c(\vX \mid \boldsymbol{\Pi})
    \doteq \sum_{k=1}^K \frac{N_k}{N}R_{\epsilon}(\vX_k)}.
\end{equation}

\paragraph{Estrategias de optimización para el agrupamiento.}
La pregunta que queda por resolver es cómo optimizar el objetivo de tasa de codificación anterior
para encontrar los grupos óptimos. Existen tres enfoques naturales
para abordar este objetivo:
\begin{enumerate}
    \item Podemos empezar con el conjunto completo $\vX$ como un único grupo
        (es decir, el régimen «perezoso») y luego buscar (por ejemplo, al azar) la forma de
        dividirlo de modo que se obtenga una tasa de codificación menor. 
    \item Por el contrario, podemos partir de cada muestra $\vx_i$ como un grupo independiente
        (es decir, el régimen de memorización)  y buscar fusionar grupos
        que den lugar a una tasa de codificación menor.
    \item Por otra parte, si pudiéramos representar (o aproximar) la
        función de pertenencia $\boldsymbol{\Pi}$ mediante unos parámetros continuos,
        podríamos utilizar métodos de optimización como el descenso por gradiente (GD).
\end{enumerate}
El primer enfoque no resulta tan atractivo desde el punto de vista computacional, ya que el número
de particiones posibles que hay que probar es exponencial con respecto al
número de muestras. Por ejemplo, el número de particiones de $\vX$ en
dos subconjuntos de igual tamaño es $N \choose N/2$, lo cual se dispara a medida que $N$
aumenta. Analizaremos el tercer enfoque en el próximo
\Cref{ch:deep-networks}. Allí veremos cómo el papel de las
redes neuronales profundas, en particular los transformadores, está relacionado con el
objetivo de la tasa de codificación.

El segundo enfoque fue propuesto originalmente en el trabajo de
\cite{ma2007segmentation}. Este trabajo demuestra las ventajas de poder
evaluar la tasa de codificación de manera eficiente (por ejemplo, mediante una fórmula
analítica). Gracias a ello, los grupos (de baja dimensión) de los datos pueden
identificarse de forma bastante eficiente y eficaz mediante el principio de
minimización de la longitud de codificación (MCL). Tenga en cuenta que, para un clúster $\vX_k$ con
$N_k$ muestras, la longitud de bits binarios necesaria para codificar todas las
muestras de $\vX_k$ viene dada por:\footnote{De hecho, una estimación más precisa
    de la longitud de codificación es $L(\vX_k) = (N_k+D) R_\epsilon(\vX_k)$
    donde los bits adicionales se utilizan para codificar la base del subespacio
\cite{ma2007segmentation}. Aquí omitimos esta sobrecarga por simplicidad.}
\begin{equation}
    L(\vX_k) = N_k R_\epsilon(\vX_k).
\end{equation}
Si tenemos dos grupos $\vX_k$ y $\vX_l$ y queremos codificar las
muestras como dos grupos separados, la longitud de bits binarios necesaria es
\begin{equation*}
    L^c(\vX_k, \vX_l) = N_k R_\epsilon(\vX_k) + N_l R_\epsilon(\vX_l) -
    N_k \log\frac{N_k}{N_k + N_l} - N_l \log\frac{N_l}{N_k + N_l}.
\end{equation*}
Los dos últimos términos corresponden al número de bits necesarios para codificar la
pertenencia de las muestras según el código de Huffman.

Entonces, dados dos grupos distintos $\vX_1$ y $\vX_2$, podemos
decidir si fusionarlos o no basándonos en la diferencia entre
las dos longitudes de codificación:
\begin{equation}
    L(\vX_k \cup \vX_l) - L^c(\vX_k, \vX_l)
\end{equation}
es positivo o negativo y $\vX_k \cup \vX_l$ denota la unión de los
conjuntos de muestras en $\bm X_k$ y $\bm X_l$. Si es positivo,
significa que la longitud de codificación se reduciría si fusionáramos los dos
grupos en uno solo. Este simple hecho da lugar al siguiente algoritmo de agrupamiento
propuesto por \cite{ma2007segmentation}:
\begin{algorithm}[!htbp]
    \caption{Pairwise Steepest Descent of Coding
    Length}\label{alg:steepest_descent_coding_length}
    \begin{algorithmic}[1]
        \Require{\(N\) puntos de datos \(\{\vx_{i}\}_{i = 1}^{N}\)}
        \Ensure{Un conjunto \(\cC\) de grupos}

        \Procedure{PairwiseSteepestDescentOfCodingLength}{$\{\vx_{i}\}_{i
        = 1}^{N}$}
        \Statex{\quad\ \texttt{\# Initializar \(N\) grupos
        \(\vX_{k}\) con un elemento cadauno}}
        \State{\(\cC \gets \{\{\bm x_{i}\}\}_{i = 1}^{N}\)}
        \Statex{}
        \While{\(\abs{\cC} > 1\)}
        \Statex{}
        \Statex{\qquad\quad \texttt{\# Si
        no hay bits ahorrados con ninguna fusión}}
        \If{\(\displaystyle \min_{\vX_{k}, \vX_{l} \in \cC}[L(\vX_{k}
        \cup \vX_{l}) -L^{c}(\vX_{k}, \vX_{l})] \geq 0\)}
        \Statex{\qquad\qquad \texttt{\# Tempranamente devuelve \(\cC\) y finaliza}}
        \State{\Return{\(\cC\)}}
        \Else
        \Statex{\qquad\qquad \texttt{\# Fusiona los grupos que
        ahorren más bits}}
        \State{\(\displaystyle \vX_{k^{\ast}}, \vX_{l^{\ast}} \gets
                \argmin_{\vX_{k}, \vX_{l} \in \cC}[L(\vX_{k} \cup \vX_{l}) -
        L^{c}(\vX_{k}, \vX_{l})]\)}
        \Statex{}
        \Statex{\qquad\qquad\texttt{\# Remover grupos no fusionados 
        y volver a agregar el grupo fusionado}}
        \State{\(\displaystyle \cC \gets [\cC \setminus
                \{\vX_{k^{\ast}}, \vX_{l^{\ast}}\}] \cup \{\vX_{k^{\ast}}
        \cup \vX_{l^{\ast}}\}\)}
        \EndIf
        \EndWhile
        \Statex{\quad\ \texttt{\# Si todas las fusiones arrojan
        ahorros, devolver un grupo}}
        \State{\Return{\(\cC\)}}
        \EndProcedure
    \end{algorithmic}
\end{algorithm}

Cabe señalar que este algoritmo es manejable, ya que el número total de
comparaciones (por pares) y fusiones es de aproximadamente $O(N^2\log N)$. Sin embargo,
debido a su naturaleza voraz, no existe garantía teórica de que el
proceso converja hacia la solución de agrupamiento óptima a nivel global.
No obstante, tal y como se indica en \cite{ma2007segmentation}, en la práctica,
este algoritmo, aparentemente sencillo, funciona extremadamente bien. Los resultados de la agrupación
representados en la \Cref{fig:two-lines-and-plane} se calcularon, de hecho,
mediante este algoritmo.

\begin{example}[Segmentación de imágenes]\label{eg:image-segmentation} La
    medida de longitud de codificación mencionada anteriormente y el algoritmo de agrupamiento
    asociado suponen que la distribución de los datos es una mezcla de
    distribuciones gaussianas (de baja dimensión). Aunque esto pueda parecer un tanto
    idealizado, la medida y el algoritmo pueden resultar muy útiles
    e incluso potentes en situaciones en las que el modelo sea (aproximadamente) válido.

    Por ejemplo, una imagen natural suele estar compuesta por múltiples
    regiones con texturas casi homogéneas. Si tomamos muchas ventanas pequeñas
    de cada región, estas deberían parecerse a muestras extraídas de
    una distribución gaussiana (de baja dimensión), tal como se ilustra en
    la \Cref{fig:image-patch}. La \Cref{fig:image-segmentation} muestra los
    resultados de la segmentación de imágenes obtenidos al aplicar el
    algoritmo de agrupamiento mencionado anteriormente directamente a los fragmentos de imagen. Se pueden encontrar más
    detalles técnicos sobre la adaptación del algoritmo al
    problema de la segmentación de imágenes en \cite{Mobahi-IJCV2011}.
\end{example}

\begin{figure}
    \centering
    \includegraphics[width=0.4\linewidth]{\toplevelprefix/chapters/chapter4/figs/image-segmentation-tiles.png}
    \caption{Fragmentos de imagen con una dimensión de $w\times w$ pixels.}
    \label{fig:image-patch}
\end{figure}

\begin{figure}[th]
    \centering
    \includegraphics[width=0.8\linewidth]{\toplevelprefix/chapters/chapter4/figs/image-segmentation.png}
    \caption{Resultados de la segmentación basados en el algoritmo de agrupamiento
    aplicado a los fragmentos de imagen.}
    \label{fig:image-segmentation}
\end{figure}

\section{Transformación mediante maximización de la ganancia de información}
\label{sec:chap4-representation-learning-problem}

Hasta ahora, en este capítulo hemos analizado cómo identificar una
distribución con estructuras de baja dimensión mediante el principio de
compresión. Como hemos visto en las dos secciones anteriores,
la compresión computacional puede llevarse a cabo mediante la
operación de eliminación de ruido o mediante la agrupación en clústeres. La \Cref{fig:Gaussian-Subspaces}
ilustra este concepto con nuestro ejemplo favorito.

Es importante señalar que el algoritmo de agrupamiento y las fórmulas de la tasa de codificación
desarrolladas en la sección anterior partían del supuesto de que la distribución de los datos
es \textit{ya} (aproximadamente) una mezcla de gaussianas de baja dimensión.
En la práctica, los datos brutos —imágenes, audio, texto— rara vez cumplen esta
hipótesis de manera directa: sus estructuras intrínsecas suelen ser
\textit{no lineales}. El objetivo de esta sección es desarrollar un
marco para \textit{transformar} los datos de tal manera que sus representaciones
adopten la forma de una mezcla de subespacios incoherentes de baja dimensión,
momento en el que las herramientas de la sección anterior se vuelven directamente
aplicables. Es decir, la sección anterior responde a la pregunta «cómo comprimir
y agrupar datos que ya son gaussianos», mientras que esta sección
responde a la pregunta «cómo transformar los datos para que se conviertan en gaussianos». Como
veremos en el \Cref{ch:deep-networks}, las redes profundas son precisamente el
mecanismo computacional que permite llevar a cabo tales transformaciones y la
arquitectura de la red se deriva de la optimización del objetivo de la tasa de codificación
desarrollado aquí.

De hecho, el objetivo final de identificar una distribución de datos es utilizarla para
facilitar ciertas tareas posteriores, como la segmentación, la clasificación o la
generación (de imágenes). Por lo tanto, la forma en que se «representa» la distribución resultante
es de vital importancia para determinar cómo se puede recuperar y
utilizar de manera eficiente y eficaz la información relacionada con estas
tareas posteriores. Esto plantea, naturalmente, una pregunta fundamental: {\em ¿qué
hace que una representación sea verdaderamente «buena» para su uso posterior?} A
continuación, analizaremos las propiedades esenciales que debe poseer una representación
significativa y útil y cómo estas propiedades
pueden caracterizarse y perseguirse explícitamente mediante la maximización de la ganancia de información.

\paragraphCómo evaluar la calidad de las representaciones.}
Se puede considerar un conjunto de datos dado como un conjunto de muestras de un vector aleatorio $\vx$ con
una distribución determinada en un espacio de alta dimensión, por ejemplo,
$\R^D$. Por lo general, la distribución de $\vx$ tiene una dimensión intrínseca mucho menor
que el espacio circundante. En términos generales,  {\em
aprender una representación} se refiere a aprender una aplicación continua,
digamos $f(\cdot)$, que transforma $\vx$ en un denominado {\em vector de características
} $\vz$ en otro espacio (normalmente de menor dimensión), digamos
$\R^d$, donde $d < D$. Se espera que a través de tal mapeo
\begin{equation}
    \vx \in \R^D \xrightarrow{\hspace{2mm} f(\vx)\hspace{2mm}}
    \vz  \in \R^d,
    \label{eqn:chap4-1-encoding}
\end{equation}
las estructuras intrínsecas de baja dimensión de $\vx$ se identifiquen y
se representen mediante $\vz$ de una manera más compacta y estructurada, con el fin de
facilitar tareas posteriores como la clasificación o la generación. La
característica $\vz$ puede considerarse como un código compacto (aprendido) de los
datos originales $\vx$, por lo que la aplicación $f$ también se denomina \textit{codificador}.
La pregunta fundamental del aprendizaje de representaciones es
\begin{center}
    \noindent{\em ¿Cuál sería una medida rigurosa y eficaz para evaluar
    la calidad de las representaciones?}
\end{center}

\begin{figure}[t]
    \centering
    \includegraphics[width=0.8\linewidth]{\toplevelprefix/chapters/chapter4/figs/Gaussian-Subspaces.png}
    \caption{Identificación de una distribución de baja dimensión con dos
        subespacios (izquierda) mediante la eliminación de ruido o la agrupación, partiendo de una
    distribución gaussiana aleatoria genérica (derecha).}
    \label{fig:Gaussian-Subspaces}
\end{figure}

Conceptualmente, la calidad de una representación $\vz$ depende de
qué tan bien identifique la información más relevante y suficiente de
$\vx$ para las tareas posteriores y qué tan eficientemente represente dicha información.
Durante mucho tiempo, se creía y se defendía que la «suficiencia»
o «calidad» de una representación de características aprendida debía definirse
en función de una tarea específica. Por ejemplo, $\vz$ solo tiene que ser
suficiente para predecir la etiqueta de clase $\y$ en un problema de clasificación.
A continuación, comenzaremos con el problema clásico de la clasificación de imágenes
y analizaremos por qué ese concepto de «representación» específica para cada tarea
es limitado y debe generalizarse.

\subsection{Aprendizaje de representaciones a partir de datos clasificados}

Supongamos que $\bm{x} \in \R^D$ es un vector aleatorio extraído de
una mezcla de $K$ distribuciones (componentes)  $\mathcal{D} =
\{\mathcal{D}_k\}_{k=1}^K$. Dado un conjunto finito de muestras independientes e idénticamente distribuidas $\vX =
[\vx_1, \vx_2, \ldots, \vx_N] \in \Re^{D\times N}$ del vector aleatorio
$\bm x$, {\em buscamos una buena representación}
mediante una aplicación continua $f(\vx): \R^D \rightarrow \R^d$ que capte
las estructuras intrínsecas de $\vx$ y facilite al máximo la posterior
tarea de clasificación.\footnote{La clasificación es el ámbito en el que el
    aprendizaje profundo demostró sus primeros éxitos, lo que desató un
    interés explosivo por las redes profundas. Aunque nuestro estudio se centra en
    la clasificación, creemos que las ideas y los principios pueden
generalizarse de forma natural a otros contextos, como la regresión.} Para facilitar la tarea
de aprender la distribución $\mathcal{D}$, en el popular contexto de la
clasificación supervisada, se asigna a cada muestra $\vx_i$ una etiqueta de clase verdadera (o una palabra clave para cada
clase), representada normalmente por un vector one-hot $\y_i \in \R^K$.

\paragraph{Codificación de la información de clases mediante la entropía cruzada.} Numerosos
estudios han demostrado que, para muchos conjuntos de datos prácticos (por ejemplo, imágenes,
audio y lenguajes naturales), el mapeo (de codificación) entre los datos
$\bm{x}$ y su etiqueta de clase $\bm{y}$ puede modelarse de manera eficaz mediante
el entrenamiento de una red profunda,\footnote{Por ahora, no nos preocupemos
    por qué red debemos usar aquí y por qué. El objetivo aquí es
    analizar cualquier red profunda probada empíricamente. Dejaremos la
justificación de las arquitecturas de red para el próximo capítulo.} aquí
denotada $$f(\vx, \theta):\vx \mapsto \y$$ con parámetros de red
$\theta \in \Theta$, donde $\Theta$ es el espacio de parámetros. Para
que la salida $f(\vx, \theta)$ coincida bien con la etiqueta $\y$, nos interesa
minimizar la {\em pérdida de entropía cruzada} sobre un conjunto de entrenamiento
$\{(\vx_i, \y_i)\}_{i=1}^N$:\footnote{Aquí interpretamos \(\vy\) como una
    distribución de probabilidad \textit{one-hot} concreta sobre \([K] =
    \{1, 2, \dots, K\}\),
    y \(f(\vx, \theta) \in \Delta([K]) \subseteq \R^{K}\) como una
    distribución de probabilidad sobre el mismo conjunto, que por lo tanto puede
    compararse mediante la entropía cruzada a pesar de estar representada computacionalmente
como vectores en \(\R^{K}\).}
\begin{equation}
    \min_{\theta \in \Theta} \; \mathbb{E}[\CE(\vy \mmid f(\vx,
    \theta))] \, \approx - \frac{1}{N}\sum_{i=1}^N \langle
    \y_i, \log\left(f(\vx_i, \theta)\right) \rangle.
    \label{chap4-eqn:cross-entropy}
\end{equation}
Los parámetros óptimos $\theta$ de la red suelen obtenerse
optimizando el objetivo anterior mediante un método eficiente de descenso de gradientes,
en el que los gradientes se calculan mediante retropropagación, tal y como
se describe en la \Cref{app:BP-section} del \Cref{app:optimization}.

A pesar de su eficacia y su enorme popularidad, este enfoque presenta dos
limitaciones importantes: (1) Su único objetivo es predecir
las etiquetas $\y$, incluso si estas pudieran estar mal etiquetadas. Los estudios empíricos
muestran que las redes profundas, utilizadas como una «caja negra», pueden incluso adaptarse a etiquetas aleatorias
~\cite{zhang2017understanding}. (2) Con este tipo de ajuste de datos de extremo a extremo,
a pesar de los numerosos esfuerzos empíricos por interpretar las características así aprendidas, %
no queda claro hasta qué punto las características intermedias aprendidas por
la red captan las estructuras intrínsecas de los datos que hacen
posible, en primer lugar, una clasificación significativa. Las propiedades geométricas y estadísticas precisas
de las características aprendidas también suelen quedar
ocultas, lo que da lugar a una falta de interpretabilidad y a la
consiguiente ausencia de garantías de rendimiento (por ejemplo, generalización,
transferibilidad y robustez, etc.) en el aprendizaje profundo.
Por lo tanto, {\em uno de los objetivos de esta sección es abordar dichas
    limitaciones reformulando el objetivo para aprender
    representaciones explícitamente significativas y útiles para los datos $\vx$,
sin limitarse a la clasificación.}

\begin{figure}
    \centering
    \includegraphics[width=0.95\linewidth]{\toplevelprefix/chapters/chapter4/figs/neural_collapse.pdf}
    \caption{Evolución de los resultados de la penúltima capa de una red neuronal VGG13
        cuando se entrena con el conjunto de datos CIFAR-10 utilizando 3 clases seleccionadas al azar.
    Figura tomada de \cite{papyan2020prevalence}.}
    \label{chap4-fig:neural-collapse}
\end{figure}

\paragraph{Características discriminativas mínimas mediante cuello de botella de información.}
Un enfoque habitual para interpretar el papel de las redes profundas consiste en
considerar que las salidas de las capas intermedias de la red seleccionan
ciertas características latentes $\vz = f(\vx, \theta) \in \Re^d$ de los datos
que permiten distinguir entre varias clases. Las representaciones
aprendidas $\vz$ facilitan entonces la tarea de clasificación
posterior para predecir la etiqueta de clase $\y$ mediante la optimización de un clasificador $g(\vz)$:
\begin{equation}
    \vx   \xrightarrow{\hspace{2mm} f(\vx, \theta)\hspace{2mm}} \vz
    \xrightarrow{\hspace{2mm} g(\vz) \hspace{2mm}} \y.
\end{equation}
Sabemos por \eqref{eq:mutual} que {\em la
información mutua} entre dos variables aleatorias, digamos $\vx,\vz$, se
define como
\begin{equation}
    I(\vx; \vz) = H(\vx) - H(\vx\mid \vz),
\end{equation}
donde $H(\vx \vert \vz)$ es la entropía condicional de $\vx$ dado $\vz$.
La información mutua también se conoce como {\em la ganancia de información}:
Mide en qué medida se puede reducir la entropía de la variable aleatoria $\vx$
una vez que se conoce $\vz$. O, dicho de otro modo, mide cuánta
información contiene $\vz$ sobre $\vx$. La formulación del {\em cuello de botella de la información
} (IB) \cite{Tishby-ITW2015} postula además
que la función de la red es aprender $\vz$ como la
estadística mínima suficiente para predecir $\y$. Formalmente, busca
maximizar la información mutua $I(\vz, \y)$
entre $\vz$ y $\y$, al tiempo que minimiza la información mutua $I(\vx,
\vz)$ entre $\vx$ y $\vz$:
\begin{equation}
    \max_{\theta\in \Theta}\; \mbox{IB}(\vx, \y, \vz) \doteq I(\vz; \y) -
    \beta I(\vx; \vz) \quad\ \mathrm{s.t.}\ \vz = f(\vx, \theta),
    \label{chap4-eqn:information-bottleneck}
\end{equation}
donde $\beta >0$.

Si se logran superar algunas limitaciones asociadas a este marco
\cite{kolchinsky2018caveats-ICLR2018}, como por ejemplo cómo evaluar con precisión
la información mutua con un número finito de muestras provenientes de distribuciones degeneradas,
este marco puede resultar útil para explicar ciertos
comportamientos de las redes profundas.
Por ejemplo, un trabajo reciente \cite{papyan2020prevalence} demuestra, de hecho,
que las representaciones aprendidas mediante la pérdida de entropía cruzada
\eqref{chap4-eqn:cross-entropy} presentan un fenómeno de \emph{colapso neuronal}.
Es decir, las características de cada clase se asignan a un vector unidimensional,
mientras que el resto de la información de la clase se omite, tal y como
se ilustra en la \Cref{chap4-fig:neural-collapse}.\footnote{Observaremos 
    este fenómeno en experimentos con datos de imágenes reales más adelante
en el \Cref{eg:Rate-Reduction-CIFAR10}.}
\begin{remark}[Neural Collapse]
    El colapso neuronal se refiere a un fenómeno observado en redes neuronales profundas
    entrenadas para la clasificación, en el que las representaciones de características
    aprendidas y los pesos del clasificador muestran un comportamiento altamente simétrico y
    estructurado durante la fase final del entrenamiento
    \cite{papyan2020prevalence,zhu2021geometric,yaras2022neural}.
    Concretamente, dentro de 
    cada clase, las características se ajustan a la media de su clase y entre
    las clases, estas medias se separan al máximo, formando una configuración
    equiangular en forma de simplex. El clasificador lineal se alinea con la
    media de la clase, salvo por un cambio de escala. Además, el clasificador de la última capa
    converge en elegir aquella clase cuya media de entrenamiento
    sea la más cercana. El colapso neuronal revela conexiones profundas entre la dinámica de la optimización,
    la generalización y las estructuras geométricas que surgen en
    el aprendizaje supervisado.
    \label{remark:neural-collapse}
\end{remark}

A partir del ejemplo de clasificación anterior, vemos que la representación
aprendida $\vz$ de $\vx$ da lugar a un codificador muy sencillo que,
en esencia, asigna a cada clase de datos una sola palabra de código: el
vector «one-hot» que representa cada clase. Desde el punto de vista de la compresión con pérdidas,
un codificador de este tipo presenta demasiadas pérdidas como para conservar la información en
la distribución de los datos. Otra información, como la que resulta útil para
tareas tales como la segmentación, la completación o la generación de imágenes, se
pierde en gran medida en un proceso de aprendizaje supervisado de este tipo.

Para remediar esta situación, queremos aprender un esquema de codificación diferente
de tal manera que la representación de características resultante $\vz$ pueda
capturar tanta información como sea posible sobre la distribución de los datos,
sin limitarse únicamente a la información útil para la clasificación. Más concretamente,
queremos que la representación $\vz = f(\vx, \theta)$ sea tal que
exista un decodificador $g(\vz,\eta)$ de modo que:
\begin{equation}
    \|\vx - \hat \vx\|_2 \le \epsilon \quad \text{s.t.} \quad \hat{\vx}
    = g(f(\vx)).
\end{equation}
Introduciremos formalmente dicho esquema de codificación-decodificación como un
«autocodificador» y dicha representación $\vz$ como una representación «consistente»
en el \Cref{ch:consistent-representations}. Por ahora, el lector
puede considerarlo simplemente como una generalización del código de cuantificación
$\hat \vx$ introducido en la sección anterior para la tasa-distorsión
(en el dominio de datos nativos $\vx$).

En la formulación del cuello de botella de información
\eqref{chap4-eqn:information-bottleneck}, no se requiere que la
representación de características aprendida $\vz$ sea consistente. Por lo tanto, como mínimo
debería fijarse un objetivo para el aprendizaje de representaciones (RL)
(consistente):
\begin{eqnarray}
    \max_{\theta,\eta}\; \mbox{RL}(\vx, \y, \vz) \doteq I(\vz; \y) - \beta
    I(\vx; \vz) \nonumber \\ \mathrm{s.t.}\quad  \vz = f(\vx,\theta),\quad
    \|\vx - g(\vz, \eta) \|_2 \le \epsilon.
    \label{chap4-eqn:representation-learning}
\end{eqnarray}
Resolver el problema anterior en el contexto más general es una tarea muy
desalentadora. En general, ni siquiera sabemos qué clases de
codificadores o decodificadores podrían garantizar la consistencia\footnote{En
    la práctica, las personas tienden a aprenderlas, al margen del objetivo
    mencionado anteriormente, mediante ensayo y error, como veremos con más detalle en el
\Cref{ch:consistent-representations}.} Además, dado que desconocemos la
distribución de la representación $\vz$, calcular la información mutua
$I(\vz;\y)$ puede resultar muy complicado, si no imposible.

Antes de analizar cómo abordar este problema en contextos generales (en
los capítulos siguientes), en el resto de este capítulo mostraremos
cómo resolverlo, de manera limitada, en el caso particular en el que
la distribución de los datos se compone de clases bien diferenciadas (de baja dimensión).

\subsection{Representaciones discriminativas lineales}\label{subsec:LDR}

Que los datos dados $\vX$ de una distribución mixta $\mathcal{D}$ puedan
clasificarse o agruparse de manera efectiva depende de cuán separables (o
discriminativas) sean (o puedan hacerse) las distribuciones componentes $\mathcal{D}_k$.
Una hipótesis de trabajo muy extendida es que la distribución
de cada clase presenta estructuras intrínsecas
de relativamente {\em baja dimensión}. Por lo tanto, podemos suponer que la distribución $\mathcal{D}_k$
de cada clase tiene su soporte en una subvariedad de baja dimensión, digamos
$\mathcal{M}_k$ con dimensión $d_k \ll D$, y que la distribución
$\mathcal D$ de $\vx$ se apoya en la mezcla de esas
subvariedades, $\mathcal M = \cup_{k=1}^K \mathcal{M}_k$,  en el
espacio ambiental de alta dimensión $\Re^D$, como se ilustró anteriormente en la
\Cref{fig:representation-learning}.

No solo necesitamos identificar la distribución de baja dimensión, sino que
también queremos representar la distribución de una forma que
facilite al máximo las tareas posteriores, como la clasificación, la agrupación en clústeres y
la generación condicionada (como veremos más adelante). Para ello,
necesitamos que nuestras representaciones de características aprendidas tengan las siguientes
propiedades:
\begin{enumerate}
    \item {\em Compresibilidad intraclase:} Las características de las muestras de la
        misma clase deben estar fuertemente {\em correlacionadas}, en el sentido de que
        pertenezcan a un subespacio lineal de baja dimensión.
    \item {\em Discriminación entre clases:} Las características de las muestras de
        diferentes clases deben estar muy {\em descorrelacionadas} y pertenecer
        a subespacios lineales de baja dimensión diferentes e incoherentes.
    \item {\em Representación con máxima diversidad:} La dimensión (o varianza)
        de las características de cada clase debe ser {\em lo más grande
        posible}, siempre y cuando sean independientes de las de las demás clases.
\end{enumerate}
Nos referimos a dicha representación como la {\emrepresentación discriminativa 
lineal} (LDR).

Obsérvese que la primera propiedad concuerda bien con el objetivo del
clásico {\em análisis de componentes principales} (PCA) que hemos
analizado en la \Cref{sub:pca}. En esencia, esto significa que las características de
las muestras que pertenecen a la misma clase\footnote{incluidas las
instancias aumentadas de las mismas muestras} deben volverse o
hacerse altamente correlacionadas. En otras palabras, la representación de características debe ser {\em
contractiva} para las muestras que son similares.

\begin{figure}
    \centering
    \includegraphics[width=0.8\linewidth]{\toplevelprefix/chapters/chapter4/figs/mcr_diagram.png}
    \caption{La distribución $\mathcal D$ de datos de alta dimensión
        $\vx\in \Re^D$ se define sobre una variedad $\mathcal{M}$ y sus
        clases se encuentran en subvariedades de baja dimensión $\mathcal{M}_k$. Nuestro objetivo es
aprender una aplicación $f(\vx, \theta)$ parametrizada por $\theta$ tal que
$\vz_i = f(\vx_i, \theta)$ se encuentre en una unión de subespacios con máxima falta de correlación
$\{\mathcal{S}_k\}$.}
    \label{chap4-fig:mcr-diagram}
\end{figure}

La segunda propiedad exige que, en el caso de las muestras que pertenecen a
clases diferentes o que se consideran distintas entre sí, sus características
deben ser (geométricamente) incoherentes o (estadísticamente)
independientes, tal y como se ilustra anteriormente en la
\Cref{fig:representation-learning}. La \Cref{chap4-fig:mcr-diagram}
ilustra  cuándo la distribución de los datos es, en realidad, una mezcla de
subvariedades de baja dimensión. Mediante la compresión (eliminación de ruido o
agrupación), identificamos en primer lugar la distribución real de los datos
(izquierda). A continuación, queremos transformar la distribución de tal manera que las
subvariedades acaben convirtiéndose en subespacios lineales mutuamente incoherentes o independientes
(derecha).

\begin{remark}[Análisis lineal discriminativo]
    Se puede observar que
    el segundo criterio se asemeja al del clásico {\em  análisis discriminante lineal
    (LDA)~\cite{HastieTiFr09}. El análisis discriminante lineal (LDA)~\cite{HastieTiFr09} es una
    técnica supervisada de reducción de dimensionalidad cuyo objetivo es encontrar una
    proyección lineal de los datos que maximice la separabilidad de las clases.
    En concreto, a partir de datos etiquetados, el LDA busca una transformación lineal
    que proyecte las entradas de alta dimensión en un espacio de menor dimensión
    en el que las clases estén separadas al máximo. Tenga en cuenta que el PCA es un
    método no supervisado que proyecta los datos en direcciones de máxima
    varianza sin tener en cuenta las etiquetas de clase. Mientras que el PCA se centra exclusivamente
    en preservar la estructura global de la varianza, el LDA aprovecha explícitamente
    la información de las etiquetas para mejorar el poder discriminativo; véase la comparación
    en la \Cref{fig:LDA}.
    \begin{figure}
        \centering
        \includegraphics[width=0.75\linewidth]{\toplevelprefix/chapters/chapter4/figs/LDA.png}\vspace{-0.1in}
        \caption{Comparación entre el PCA y el LDA. Figuras tomadas de
        \url{https://sebastianraschka.com/Articles/2014_python_lda.html}.}
        \label{fig:LDA}
    \end{figure}
    Pero el LDA solo realiza una proyección lineal de los datos originales. Como
    veremos, el LDR puede ser más general: busca una transformación, probablemente no lineal,
    para que la representación resultante sea lineal
    e incoherente.
\end{remark}

\begin{remark}[Aprendizaje contrastivo]
    Cabe señalar que el {\emaprendizaje contrastivo} convencional
    \citep{hadsell2006dimensionality,Van_den_Oord2018-qf,he2019momentum}
    pretende alcanzar objetivos similares en cuanto a la representación de características, tal y como establecen las dos primeras
    propiedades. En el caso de datos que pertenecen a $k$ clases diferentes,
    un par elegido al azar $(\vx_i, \vx_j)$ tiene una alta probabilidad
    de pertenecer a clases diferentes si $k$ es grande.\footnote{Por
        ejemplo, cuando $k \ge 100$, un par aleatorio tiene un 99\% de probabilidad
    de pertenecer a clases diferentes.} Por lo tanto, es deseable que la
    representación $\vz_i = f(\vx_i, \bm \theta)$ de una muestra $\vx_i$ sea
    altamente incoherente con respecto a las $\vz_j$ de otras muestras $\vx_j$, mientras que
    coherente con la característica de su versión transformada $\tau(\vx_i)$, denotada
    como $\vz(\tau(\vx_i))$ para $\tau$ en cierto conjunto de aumento $\mathcal
    T$ en consideración. Por lo tanto, se propuso de manera heurística que, para
    fomentar el poder discriminativo de la representación aprendida, se podría
    intentar minimizar la denominada {\em pérdida contrastiva}:
    \begin{equation}
        \min_{\bm \theta} - \log \frac{\exp(\langle \vz_i,
        \vz(\tau(\vx_i))\rangle)}{\sum_{j\not= i} \exp(\langle \vz_i,
        \vz_j\rangle)},
        \label{eqn:contrastive-loss}
    \end{equation}
    que es pequeña siempre que el producto interno $\langle \vz_i,
    \vz(\tau(\vx_i))\rangle$ sea grande y $\langle \vz_i, \vz_j\rangle$ sea
    pequeño para $i \neq j$. Esta pérdida contrastiva y muchas de sus
    variantes se han utilizado ampliamente en métodos populares de aprendizaje de características
    (autosupervisado), como las características de tipo DINO
    \cite{caron2021unsupervisedlearningvisualfeatures,DINOv2} para datos visuales.
    Como veremos, basándonos en el marco de codificación con pérdida estudiado en
    la sección anterior, podremos llegar a una función objetivo fundamentada
    que promueva ambas propiedades y, al mismo
    tiempo, simplifique significativamente la práctica, como el aprendizaje de
    las características de tipo DINO (tal y como expondremos con gran detalle en el
    \Cref{ch:applications}).
\end{remark}

La tercera propiedad también es importante porque queremos que las características aprendidas
revelen todas las causas posibles por las que una clase es diferente
de todas las demás. Por ejemplo, para distinguir una «manzana» de
una «naranja», no solo nos fijamos en el color, sino también en la forma y en las
hojas. Lo ideal sería que la dimensión de cada subespacio $\{\mathcal{S}_k\}$
fuera igual a la de la subvariedad correspondiente
$\mathcal{M}_k$. Esta propiedad será importante si queremos que la
función $f(\vx,\theta)$ sea {\em invertible} para tareas como la
generación de imágenes. Por ejemplo, si tomamos diferentes puntos de muestra del
subespacio de características de «manzana», deberíamos poder decodificarlos, mediante
una aplicación inversa $g(\vz, \eta)$, para generar diversas imágenes de manzanas.
Este par de aplicaciones mutuamente inversas $f$ y $g$ constituye un
autocodificador que estudiaremos con gran detalle en el
\Cref{ch:consistent-representations}. Obsérvese que las representaciones de características
aprendidas a partir de la minimización de la entropía cruzada \eqref{chap4-eqn:cross-entropy}
normalmente no presentan esta propiedad, ya que adolecen del denominado fenómeno del {\em
colapso neuronal} comentado anteriormente (véase la
\Cref{remark:neural-collapse}).

En general, aunque las estructuras intrínsecas de cada clase o grupo
puedan ser de baja dimensión, no son en absoluto simplemente lineales (o
gaussianas) en su representación original $\vx$, y es necesario
convertirlas primero en lineales mediante alguna transformación no lineal.\footnote{Analizaremos
    cómo se puede hacer esto de manera explícita en el
\Cref{ch:consistent-representations}.} Por lo tanto, en general, utilizamos
la transformación no lineal
$f(\vx,\theta)$ para buscar una representación de los datos
tal que los subespacios que representan todas las clases sean subespacios lineales
máximamente incoherentes. Para ser más precisos, queremos aprender una
función {$\vz = f(\vx,\theta)$} que mapee cada una de las subvariedades
$\mathcal{M}_k \subset \Re^D$ (\Cref{chap4-fig:mcr-diagram} izquierda) a
un subespacio {\em lineal} $\mathcal{S}_k \subset \Re^d$
(\Cref{chap4-fig:mcr-diagram} derecha). En cierta medida, los
subespacios múltiples resultantes $\{\mathcal{S}_k\}$ pueden considerarse como
{\em componentes principales generalizadas} discriminativas \cite{GPCA} o,
si son ortogonales, {\em componentes independientes}
\cite{hyvarinen2000independent} de las características resultantes $\vz$ para
los datos originales $\bm x$.
Como veremos en el siguiente \Cref{ch:deep-networks}, las redes profundas
desempeñan precisamente la función de modelar y llevar a cabo esta transformación no lineal
de la distribución de los datos a representaciones discriminativas lineales.

\subsection{El principio de la reducción máxima de la tasa de 
codificación}\label{subsec:MCR2}

Aunque las tres propiedades —{\em discriminación entre clases}, {\em
compresibilidad dentro de las clases} y {\em representación de máxima diversidad
}---de las representaciones discriminativas lineales (LDR)
son todas propiedades muy deseables de la representación aprendida
$\vz$, no son en absoluto fáciles de obtener: ¿Son estas propiedades
compatibles entre sí, de modo que podamos esperar alcanzarlas todas a la vez? Si es así,
¿existe algún objetivo {\em sencillo pero basado en principios} que permita evaluar la
calidad de las representaciones resultantes en función de todas estas
propiedades? La clave para responder a estas preguntas {es encontrar} una
«medida de compacidad» o «ganancia de información» fundamentada para la
distribución de una variable aleatoria $\vz$ o a partir de sus muestras finitas
$\{\bm z_i\}_{i=1}^N$. Dicha medida debería caracterizar de manera directa y precisa
las propiedades geométricas o estadísticas intrínsecas de la
distribución, en términos de su dimensión intrínseca o {volumen}. A diferencia
de la entropía cruzada \eqref{chap4-eqn:cross-entropy} o el cuello de botella de información
\eqref{chap4-eqn:information-bottleneck}, dicha medida
no debe depender exclusivamente de las etiquetas de clase, de modo que pueda funcionar en
entornos más generales, como los supervisados, autosupervisados,
semisupervisados y no supervisados (como veremos en
secciones posteriores).

Sin pérdida de generalidad, supongamos que la distribución $\mathcal D$
del vector aleatorio $\vx$ se define sobre una mezcla de distribuciones,
es decir, $\mathcal D = \cup_{k=1}^K \mathcal{D}_k$, donde cada
$\mathcal{D}_k \subset \Re^D$ tiene una dimensión intrínseca baja en el
espacio envolvente de alta dimensión $\Re^D$. Sea $\bm X_k \in \Re^{D\times
N_k}$ la matriz de datos cuyas columnas son muestras extraídas de la
distribución $\mathcal{D}_k$, donde $N_k$ indica el número de
muestras para cada $k=1,\dots,K$. Entonces, utilizamos $\bm X=[\bm
X_1,\dots,\bm X_K] \in \Re^{D\times N}$ para denotar todas las muestras,
donde $N=\sum_{k=1}^K N_k$. 
Recordemos que también utilizamos $\vx_i$ para denotar la $i$-ésima muestra de $\vX$,
es decir, $\vX=[\vx_1,\dots,\vx_N] $. Bajo una aplicación de codificación:
\begin{equation}
    \vx  \xrightarrow{\hspace{2mm} f(\vx)\hspace{2mm}} \vz,
\end{equation}
Las muestras de entrada se mapean a $\vz_i = f(\vx_i)$ para cada
$i=1,\dots,N$. Abusando un poco de la notación, también escribimos $\vZ_k =
f(\vX_k)$ y $\vZ = f(\vX)$. Por lo tanto, tenemos $\vZ = [\bm
Z_1,\dots,\bm Z_K]$ y $\vZ = [\vz_1,\dots\vz_N]$.

Por un lado, para que las características aprendidas sean discriminativas, es preferible que las características de
las diferentes clases o grupos sean {\em lo más
incoherentes posible} entre sí. Por lo tanto, en conjunto deberían abarcar un espacio
del mayor volumen (o dimensión) posible y la tasa de codificación
de todo el conjunto $\vZ$ debería ser lo más alta posible. Por otro lado,
las características aprendidas de una misma clase o clúster deberían estar muy
correlacionadas y ser coherentes. Por lo tanto, cada clase o grupo solo debe abarcar un
espacio (o subespacio) de volumen muy reducido y la tasa de codificación debe
ser lo más baja posible. A continuación, explicaremos cómo medir la
tasa de codificación de las características aprendidas.

\paragraph{La tasa de codificación de las características.} Cabe destacar que una dificultad práctica
a la hora de evaluar la tasa de codificación es que la distribución subyacente de
las representaciones de las características $\vZ$ suele ser desconocida. Para solucionar
esto, podemos aproximar las características $\vZ = [\vz_1, \ldots, \vz_N]$ como
muestras extraídas de una distribución gaussiana multivariante. Bajo esta
hipótesis, tal y como se analiza en la \Cref{subsec:lossy DR}, la compacidad
de las características $\vZ$ {\em en su conjunto} puede medirse en términos de la
longitud media de codificación por muestra, denominada {\em tasa de codificación}, 
sujeta a un nivel de precisión $\epsilon > 0$ (véase
\eqref{eqn:rate-Gaussian}) definido de la siguiente manera:
\begin{equation}
    R_{\epsilon}(\vZ) = \frac{1}{2}\logdet\left(\I +
    \frac{d}{N\epsilon^{2}}\vZ\vZ^{\top}\right).
    \label{chap4-eqn:coding-length-eval}
\end{equation}

Por otro lado, esperamos que una transformación no lineal $f(\vx)$
asigne a cada subvariedad específica de clase $\mathcal{M}_k \subset
\R^D$ un subespacio lineal maximamente incoherente
$\mathcal{S}_k \subset \R^D$, de tal manera que las características aprendidas
$\vZ$ se encuentren en una unión de subespacios de baja dimensión. Esta estructura
permite una evaluación más precisa de la tasa de codificación mediante el análisis
separado de cada subespacio. %
Recordemos que las columnas de $\vZ_k$ representan las características de las muestras
en $\vX_k$ para cada $k=1,\dots,K$. La tasa de codificación de las características en
$\bm Z_k$ se puede calcular de la siguiente manera:
\begin{align}
    R_{\epsilon}(\vZ_k) = \frac{N_k}{2N}\log\det\left(\I +
    \frac{d}{N_k\epsilon^{2}}\vZ_k\vZ_k^{\top}\right)
\end{align}
Entonces, la suma de las tasas de codificación promedio de las características de cada clase es
\begin{equation}
    R_{\epsilon}^c(\vZ) \doteq \sum_{k=1}^K R_{\epsilon}(\vZ_k),
    \label{chap4-eqn:compress-loss-eval}
\end{equation}

Por lo tanto, una buena representación $\vZ$ de $\vX$ es aquella que
logra una gran diferencia entre la tasa de codificación del conjunto y
la de todas las clases:
\begin{equation}
    \Delta R_{\epsilon}(\vZ) \doteq R_{\epsilon}(\vZ) - R_{\epsilon}^c(\vZ).
    \label{chap4-eqn:coding-length-reduction}
\end{equation}
Obsérvese que, tal y como hemos comentado anteriormente en este capítulo, esta
diferencia puede interpretarse como la cantidad de «información ganada»
al identificar los clústeres de baja dimensión correctos $\vZ_k$ dentro del
conjunto global $\vZ$.

Si elegimos que nuestra función de mapeo de características $f(\cdot)$ sea una red neuronal profunda
$f(\cdot,\theta)$ con parámetros de red $\theta$, el
proceso general de representación de características y la consiguiente reducción de la tasa
se pueden ilustrar mediante el siguiente diagrama:
\begin{equation}
    \vX
    \xrightarrow{\hspace{2mm} f(\vx, \theta)\hspace{2mm}} \vZ
    \xrightarrow{\hspace{2mm} \epsilon \hspace{2mm}} \Delta
    R_{\epsilon}(\vZ).
    \label{chap4-eqn:flow}
\end{equation}
Tenga en cuenta que $\Delta R_{\epsilon}$ es {\em monótono} en la escala de
las características $\vZ$. Para garantizar una comparación justa entre las diferentes
representaciones, es esencial {\em normalizar la escala} de las
características aprendidas. Esto se puede lograr ya sea imponiendo que la
norma de Frobenius de cada clase $\vZ_k$ sea proporcional al número de
características en $\vZ_k \in \mathbb R^{d \times N_k}$, es decir, $\|\vZ_k\|_F^2
= N_k$, o bien normalizando cada característica para que se encuentre en la esfera unitaria,
es decir, $\vz_i \in \mathbb{S}^{d-1}$, donde $N_k$
denota el número de muestras en la clase $k$-ésima. Esta formulación
ofrece una justificación natural de la necesidad de la «normalización por lotes»
en la práctica del
entrenamiento de redes neuronales profundas \cite{ioffe2015batch}. %

Una vez que las representaciones son comparables, el objetivo pasa a ser aprender un
conjunto de características $\vZ  = f(\vX, \theta)$  tal que maximicen la
reducción entre la tasa de codificación de todas las características y la de la suma
de características con respecto a sus clases:
\begin{equation}
    \begin{aligned}
        \max_{\theta } & \;  \Delta R_{\epsilon}\big(\vZ \big) \doteq
        R_{\epsilon}(\vZ) - R_{\epsilon}^c(\vZ ), \\
        \mbox{s.t.} & \ \ \, \vZ = f(\bm X, \theta),\  \|\vZ_k\|_F^2
        = N_k,\ k=1,\dots,K.
    \end{aligned}
    \label{eqn:maximal-rate-reduction}
\end{equation}
Nos referimos a esto como el principio de la {\emreducción máxima de la tasa de codificación
} (MCR$^2$), como una encarnación de la famosa cita de Aristóteles:
\begin{quote}
    \centering
    ``{\em El todo es más que la suma de sus partes.}''
\end{quote}
Para aprender la mejor representación, se requiere que {\em el todo sea
máximamente mayor que la suma de sus partes}. En gran medida, la
cantidad $\Delta R$ puede considerarse como una realización computable del
término de información mutua $I(\vz; \y)$ introducido en el objetivo
\eqref{chap4-eqn:representation-learning}.

\begin{example}
    Para aclarar mejor estas ideas, analicemos el ejemplo que se muestra en la
    \Cref{fig:informative-representation}.
    \begin{figure}[t]
        \centering
        \includegraphics[width=0.99\linewidth]{\toplevelprefix/chapters/chapter4/figs/Expansion2.png}
        \caption{Tras identificar la distribución de los datos en baja dimensión,
            nos gustaría transformar aún más la distribución de los datos en una
            representación estructurada más informativa: $R$ es el número de
            bolas de radio $\epsilon$ que cubren todo el espacio y $R^c$
            es la suma de
            los números correspondientes a todos los subespacios (las bolas verdes). $\Delta R$
            es su diferencia (el número de bolas azules).
        }\label{fig:sphere-packing}
        \label{fig:informative-representation}
    \end{figure}
    Desde el punto de vista de la compresión, la representación de la derecha es
    {\em la más compacta} en el sentido de que la diferencia entre
    la tasa de codificación cuando todas las características se codifican como una sola gaussiana
    (azul) y la que se obtiene cuando las características se agrupan adecuadamente y se codifican
    como dos subespacios separados (verde) es máxima.\footnote{Intuitivamente,
        el cuociente entre el «volumen» de todo el espacio abarcado por todas las
    características y el que realmente ocupan las características es máximo.}
\end{example}

Tenga en cuenta que el principio MCR$^2$ mencionado anteriormente está diseñado para problemas de aprendizaje supervisado,
en los que se conocen las pertenencias a grupos (o etiquetas de clase).
Existen diversas formas de aplicar este principio, o generalizaciones adecuadas,
de manera que se eluda este requisito de las etiquetas de clase;
analizaremos estas opciones más adelante en este capítulo y en el siguiente
\Cref{ch:deep-networks}. Para tener una primera idea, extendemos el
principio a problemas de aprendizaje no supervisado mediante la introducción de una
matriz de pertenencia, que codifica la asignación (potencialmente flexible) de
cada punto de datos a grupos o clústeres latentes. Concretamente, sea $\bm
\Pi = \{\bm \Pi_k\}_{k=1}^K \subset \R^{N\times N}$ un conjunto de
matrices diagonales cuyos elementos diagonales codifican la pertenencia de las
$N$ muestras a $K$ clases. Es decir, $\bm \Pi$ se encuentra en un simplex
$\Omega \doteq \{\bm \Pi: \bm \Pi_k \ge \bm 0: \sum_{k=1}^K \bm \Pi_k
= \bm I_N\}$. Entonces, podemos definir la tasa de codificación promedio con
respecto a la partición $\bm \Pi$ como
\begin{align}\label{eq:MCRc}
    R_{\epsilon}^c(\vZ \mid \bm \Pi) \doteq \sum_{k=1}^K
    \frac{\mathrm{tr}(\bm \Pi_k)}{2N}\log\det\left(\bm I +
        \frac{d}{\mathrm{tr}(\bm \Pi_k)\epsilon^2}\bm Z\bm \Pi_k \bm
    Z^\top \right).
\end{align}
Cuando se da $\bm Z$, $R_{\epsilon}^c(\vZ \vert \bm \Pi)$ es una
función cóncava de $\bm \Pi$. Entonces, el principio MCR$^2$ para
problemas de aprendizaje no supervisado queda así:
\begin{align}\label{eq:MCR pi}
    \max_{\bm \Pi, \theta} & \  \Delta R_{\epsilon}\big(\vZ  \mid \bm \Pi)
    \doteq R_{\epsilon}(\bm Z) - R_{\epsilon}^c(\vZ \mid \bm \Pi) \notag \\
    \mathrm{s.t.}  & \ \ \ \bm Z = f(\bm X, \theta),\ \|\bm Z \bm
    \Pi_k\|_F^2 = N_k,\ k = 1,\dots,K,\ \bm\Pi \in \Omega.
\end{align}
En comparación con \eqref{eqn:maximal-rate-reduction}, la formulación aquí
permite la optimización conjunta tanto de las pertenencias a grupos como
de los parámetros de la red. En particular, cuando $\bm \Pi$ se fija como una
matriz de pertenencia a grupos que asigna $N$ puntos de datos a $K$ grupos,
el problema \eqref{eq:MCR pi} puede recuperar el problema
\eqref{eqn:maximal-rate-reduction}.

\begin{figure}[t]
    \centering
    \includegraphics[width=0.9\linewidth]{\toplevelprefix/chapters/chapter4/figs/mcr2-global.png}
    \caption{{\bf Landscape de optimización local:} Según
        \Cref{thm:MCR2-properties}, el máximo global del objetivo de reducción de la tasa
    corresponde a una solución con subespacios mutuamente incoherentes.}
    \label{fig:mcr-global}
\end{figure}

\begin{remark}[El cálculo de las tasas de codficación]
    Es un hecho bien conocido que calcular el logaritmo del determinante de una
    matriz \(n \times n\) requiere un tiempo de orden \(\cO(n^{3})\). Por lo tanto,
    calcular las tasas de codificación \(\Delta R_{\epsilon}\) o
    \(R_{\epsilon}\) (etc.) parece una tarea inmanejable a la escala del análisis de datos
    moderno. Hay varias formas de solucionar esto, que
    utilizaremos sin más comentarios en los experimentos que sirven de base al resto del libro.
    \begin{itemize}
        \item \textit{Utilizar la estructura de la tasa de codificación.} Supongamos
            que \(\vZ\) es un elemento de \(\R^{d \times n}\). Entonces se cumple que
            \begin{equation}
                \logdet(\vI + \alpha \vZ\vZ^{\top}) = \logdet(\vI +
                \alpha \vZ^{\top}\vZ)
            \end{equation}
            por lo que calcular el producto matricial y las tasas de codificación
            requiere un tiempo del orden de \(\min\{d^{2}n +
            d^{3}, dn^{2} + n^{3}\}\).
        \item \textit{Utilizar la simetría del Gramiano.}
            Obsérvese que \(\vZ\vZ^{\top}\) (y, de manera similar,
            \(\vZ^{\top}\vZ\)) es una matriz simétrica. Por lo tanto, resulta
            relativamente sencillo (aunque sigue siendo de tiempo cúbico) calcular sus
            valores propios \(\lambda_{i}(\vZ\vZ^{\top})\). Entonces, el
            logaritmo del determinante se puede expresar de la siguiente manera:
            \begin{equation}
                \logdet(\vI + \alpha \vZ\vZ^{\top}) = \sum_{i =
                1}^{d}\log(1 + \alpha \lambda_{i}(\vZ\vZ^{\top})),
            \end{equation}
            (y análogmente para \(\vZ^{\top}\vZ\)). Sin embargo, este
            método queda obsoleto ante el siguiente, que aprovecha
            mejor la estructura de los datos de entrada y ofrece mejores resultados.
        \item \textit{Utilizar la positividad definida de los términos que intervienen en el
            determinante logarítmico.} Dado que \(\vZ\vZ^{\top}\) (y, de manera análoga,
            \(\vZ^{\top}\vZ\)) es simétrico y semidefinido positivo,
            la cantidad \(\vI + \alpha \vZ\vZ^{\top}\) es simétrica
            y definida positiva. Esto significa que podemos calcular la descomposición de Cholesky
            \(\vL\vL^{\top} = \vI + \alpha
            \vZ\vZ^{\top}\), donde \(\vL \in \R^{d \times d}\) es una
            matriz triangular inferior con entradas diagonales positivas. Entonces
            \begin{equation}
                \logdet(\vI + \alpha \vZ\vZ^{\top}) = 2\sum_{i =
                1}^{d}\log(L_{ii}).
            \end{equation}
            En la práctica, este método suele ser el más económico,
            especialmente cuando también se decide estratégicamente si
            calcular \(\vZ\vZ^{\top}\) o \(\vZ^{\top}\vZ\), aunque
            sigue teniendo una complejidad temporal cúbica.
    \end{itemize}
\end{remark}

\vspace{-0.1in}

\subsection{Optimization Properties of Coding Rate Reduction}
\label{sec:MCR-landscape}

In this subsection, we study the optimization properties of the
MCR$^2$ function by analyzing its optimal solutions and the structure
of its optimization landscape. To get around the technical difficulty
introduced by the neural networks, we consider a simplified version
of Problem \eqref{eqn:maximal-rate-reduction} as follows:
\begin{align}\label{eqn1:maximal-rate-reduction}
    \max_{\bm Z}\ R_{\epsilon}(\vZ) - R_{\epsilon}^c(\vZ)\qquad
    \mathrm{s.t.}\quad \|\bm Z_k\|_F^2 = N_k,\ k =1,\dots,K.
\end{align}
In theory, the MCR$^2$ principle \eqref{eqn1:maximal-rate-reduction}
benefits from great generalizability and can be applied to
representations $\vZ$ of {\em any} distributions  as long as the rates
$R_\epsilon$ and $R^c_\epsilon$ for the distributions can be
accurately and efficiently evaluated. The optimal representation
$\vZ^{\ast}$  should have some interesting geometric and statistical
properties. We here reveal nice properties of the optimal
representation with the special case of subspaces, which have many
important use cases in machine learning. When the desired
representation for $\vZ$ is multiple subspaces, the rates $R_\epsilon$
and $R^c_\epsilon$ in \eqref{eqn1:maximal-rate-reduction} are given
by \eqref{chap4-eqn:coding-length-eval}
and \eqref{chap4-eqn:compress-loss-eval}, respectively. At the
maximal rate reduction, MCR$^2$ achieves its optimal representations,
denoted as $\vZ^{\ast} = [\vZ_1^*,\dots,\vZ_K^*]$ with
$\rank{(\vZ_{k}^*)}\le d_k$. One can show that $\vZ^{\ast}$ has the
following desired properties (see \cite{yu2020learning} for a formal
statement and detailed proofs).

\begin{theorem}[\bf Characterization of Global Optimal Solutions]
    Suppose $\vZ^{\ast} = [\vZ_1^*,\dots,\vZ_K^*]$ is a global optimal
    solution of Problem~\eqref{eqn1:maximal-rate-reduction}. The
    following statements hold:
    \begin{itemize}
        \item {\em Between-Class Discriminative}: As long as the ambient
            space is adequately large ($d \ge \sum_{k=1}^{K} d_k$), the
            subspaces are all orthogonal to each other, {\em i.e.},
            $(\vZ_{k}^*)^{\top} \vZ_{l}^* = \bm{0}$ for $k \not= l$.
        \item {\em Maximally Diverse Representation}:
            As long as the coding precision is adequately high, i.e.,
            $\epsilon ^4 < c\cdot \min_{k}\left\{
            \frac{N_k}{N}\frac{d^2}{d_k^2}\right\}$, where $c>0$ is a
            constant. Each subspace achieves its maximal dimension, i.e.
            $\mathrm{rank}{(\vZ_{k}^*)}= d_k$. In addition, the largest
            $d_k-1$ singular values of $\vZ_{k}^*$ are equal.
            \label{thm:MCR2-properties}
    \end{itemize}
\end{theorem}

This theorem indicates that the MCR$^2$ principle promotes embedding
of data into multiple independent subspaces (as illustrated in
\Cref{fig:mcr-global}), with features distributed {\em isotropically}
in each subspace (except for possibly one dimension). Notably, this
theorem also confirms that the features learned by the MCR$^2$
principle exhibit the desired low-dimensional discriminative
properties discussed in \Cref{subsec:LDR}. In addition, among all
such discriminative representations, it prefers the one with the
highest dimensions in the ambient space. This is substantially
different from the objective of information
bottleneck~\eqref{chap4-eqn:information-bottleneck}.

\paragraph{Regularized MCR$^2$.}
The above theorem characterizes properties of the global optima of
the rate reduction objectives. What about other optima, such as local
ones?  Due to the constraints of the Frobenius norm, it is a
difficult task to analyze Problem \eqref{eqn1:maximal-rate-reduction}
from an optimization-theoretic perspective. Therefore, we consider
the Lagrangian formulation of \eqref{eqn1:maximal-rate-reduction}.
This can be viewed as a tight relaxation or even an equivalent
problem of \eqref{eqn1:maximal-rate-reduction} whose optimal
solutions agree under specific settings of the regularization
parameter; see \cite[Proposition 3.2]{wang2024global}.
Specifically, the formulation we study, referred to henceforth as the
\textit{regularized MCR$^2$ problem}, is as follows:
\begin{align}\label{eq:MCR-reg}
    \max_{\bm Z}\ R_{\epsilon}(\vZ) - R_{\epsilon}^c(\vZ) -
    \frac{\lambda}{2}\|\bm Z\|_F^2,
\end{align}
where $\lambda > 0$ is the regularization parameter. Although the
program  \eqref{eq:MCR-reg} is highly nonconcave and involves matrix
inverses in its gradient computation, we can still explicitly
characterize its local and global optima as follows (see
\cite[Theorem 3.1]{wang2024global} for detailed proofs).

\begin{theorem}[\bf Local and Global Optima]\label{thm:mcr-global-opt}
    Let $N_k$ denote the number of training samples in the $k$-th class
    for each $k \in \{1,\dots,K\}$, $N_{\max} \doteq
    \max\{N_1,\dots,N_K\}$, $\alpha=d/(N\epsilon^2)$, and $\alpha_{k} =
    d/(N_k\epsilon^2)$ for each $k \in \{1,\dots,K\}$. Given a coding
    precision $\epsilon > 0$, if the regularization parameter satisfies
    \begin{align}\label{eq:lambda}
        \lambda \in \left(0,
            \frac{d(\sqrt{N/N_{\max}}-1)}{N(\sqrt{N/N_{\max}}+1)\epsilon^2}
        \right],
    \end{align}
    then the following statements hold: \\
    (i) ({\bf Local maximizers}) $\bm Z^* = \left[\bm Z_1^*,\dots,\bm
    Z_K^* \right]$ is a local maximizer of Problem \eqref{eq:MCR-reg} if
    and only if the $k$-th block admits the following decomposition
    \begin{align}\label{eq:Zk opti}
        \bm Z_k^* = \left(\frac{ \eta_k + \sqrt{\eta_k^2 -
        4\lambda^2N/N_k}}{2\lambda \alpha_{k}}\right)^{1/2} \bm
        U_k \bm V_k^\top,
    \end{align}
    where (a) $r_k = \mathrm{rank}(\bm Z_k^*)$ satisfies $r_k \in
    [0,\min\{N_k,d\})$ and $\sum_{k=1}^K r_k \le \min\{N,d\}$, (b) $\bm
    U_k \in \mathcal{O}^{d \times r_k}$ satisfies $\bm U_k^{\top}\bm U_l
    = \bm 0$ for all $k \neq l$, $\bm V_k \in \mathcal{O}^{N_k \times
    r_k}$, and (c) $\eta_k := (\alpha_k-\alpha) - \lambda(N/N_k+1)$ for each
    $k\in \{1,\dots,K\}$.
    \\
    (ii) ({\bf Global maximizers}) $\bm Z^* = \left[\bm Z_1^*,\dots,\bm
    Z_K^* \right]$ is a global maximizer of Problem \eqref{eq:MCR-reg} if
    and only if (a) it satisfies all the above conditions and
    $\sum_{k=1}^K r_k = \min\{N,d\}$, and (b) for all $k \neq l \in [K]$
    satisfying $N_k < N_l$ and $r_l > 0$, we have $r_k = \min\{N_k,d\}$.
\end{theorem}

\begin{figure}[t]
    \centering
    \includegraphics[width=0.9\linewidth]{\toplevelprefix/chapters/chapter4/figs/mcr2-global-local.png}
    \caption{{\bf Global optimization landscape:} According to
        \cite{sun2015nonconvex,lee2016gradient},
        \Cref{thm:mcr-global-opt,thm:mcr-benign-opt-landscape}, both global
        and local maxima of the (regularized) rate reduction objective
        correspond to a solution with mutually incoherent subspaces. All
    other critical points are strict saddle points.}
    \label{fig:mcr-global-local}
\end{figure}

This theorem explicitly characterizes the local and global optima of
problem \eqref{eq:MCR-reg}. Intuitively, this shows that the features
represented by each local maximizer of Problem \eqref{eq:MCR-reg} are
low-dimensional and discriminative. Although we have characterized
the local and global optimal solutions in
\Cref{thm:mcr-global-opt}, it remains unknown whether these
solutions can be efficiently computed using GD to solve the problem
\eqref{eq:MCR-reg}, since GD may get stuck at other critical points
such as a saddle point.
Fortunately, \cite{sun2015nonconvex,lee2016gradient} showed that if a
function is twice continuously differentiable and satisfies the {\em
strict saddle property}, i.e., each critical point is either a local
maximizer or a strict saddle point\footnote{We say that a critical
    point is a strict saddle point of Problem \eqref{eq:MCR-reg} if it
    has a direction with strictly positive
    curvature~\cite{sun2015nonconvex}. This includes classical saddle
    points with strictly positive curvature as well as strict local
minimizers.}, GD converges to a local maximizer almost surely with
random initialization. We investigate the global optimization
landscape of the problem \eqref{eq:MCR-reg} by characterizing all its
critical points as follows (see \cite[Theorem 3.3]{wang2024global}
for detailed proofs).

\begin{theorem}[\bf Benign Global Optimization
    Landscape]\label{thm:mcr-benign-opt-landscape}
    Given a coding precision $\epsilon > 0$, if the regularization
    parameter satisfies \eqref{eq:lambda},
    it holds that any critical point $\bm Z$ of the problem
    \eqref{eq:MCR-reg} is either a local maximizer or a strict saddle point.
\end{theorem}
Together, the above two theorems
show that the learned features associated with each local
maximizer of the rate
reduction objective --- not just global maximizers --- are structured
as incoherent low-dimensional subspaces. Furthermore, the
(regularized) rate reduction objective
\eqref{eqn:maximal-rate-reduction} has a very benign landscape with
only local maxima and strict saddles as critical points, as
illustrated in \Cref{fig:mcr-global-local}.
According to \cite{sun2015nonconvex,lee2016gradient},
\Cref{thm:mcr-global-opt,thm:mcr-benign-opt-landscape} imply that
low-dimensional and discriminative representations (LDRs) can be
efficiently found by applying (stochastic) gradient descent to the
rate reduction objective \eqref{eqn:maximal-rate-reduction} from
random initialization. These results also indirectly explain why in
\Cref{eg:Rate-Reduction-CIFAR10}, if the chosen network is expressive
enough and trained well, the resulting representation typically gives
an incoherent linear representation that likely corresponds to the
globally optimal solution.
Interested readers are referred to \cite{wang2024global} for proofs.

\section{Learning Representations for Imagery Data}
\label{sec:representations-image}
\subsection{Supervised Representation Learning}
\label{sec:MCR2-classification}
\label{eg:Rate-Reduction-CIFAR10}

We here present how the MCR$^2$ objective
\eqref{eqn:maximal-rate-reduction} helps learn better representations
than the cross entropy (CE) \eqref{chap4-eqn:cross-entropy} for image
classification. Here, for simplicity,\footnote{We will rigorously
    study and justify how design or derive a proper deep architecture
from and for the MCR$^2$ objective in \Cref{ch:deep-networks}.}
we adopt the popular neural network architecture, the
ResNet-18~\cite{he2016deep}, to model the feature mapping $\vz =
f(\vx,\theta)$. We optimize the neural network parameters $\theta$ via
the backpropagation algorithm (detailed in
\Cref{app:BP-section}) to maximize the coding rate reduction. More
implementation details of this experiment can be found in
\cite{chan2021redunet}.

\Cref{fig:train-test-loss} illustrates how the two rates and their
difference (for both training and test data) evolve over epochs of
training: After an initial phase, $R_\epsilon$ gradually increases
while $R^c_\epsilon$ decreases, indicating that features $\bm Z$ are
expanding as a whole while each class $\bm Z_k$ is being compressed.

We first evaluate the performance with the CIFAR-10 image
classification dataset~\cite{krizhevsky2009learning} and the results
are shown in the table of \Cref{fig:CIFAR10-performance}. As we
can see, the MCR$^2$ objective achieves a comparable performance as
the cross entropy (CE) loss. In addition, the MCR$^2$ objective is
more robust to corruptions in the class labels since it learns a
joint (subspace) representation for each class instead of trying to
fit the labels of individual samples.

\begin{figure}
    \centering
    \includegraphics[width=0.6\textwidth]{\toplevelprefix/chapters/chapter4/figs/loss_log.pdf}

    \caption{Evolution of $R_\epsilon$, $R^c_\epsilon$, $\Delta
    R_\epsilon$ during the training process.}
    \label{fig:train-test-loss}
\end{figure}

\begin{figure}
    \centering
    \includegraphics[width=\linewidth]{chapters/chapter4/figs/Table-label-corruption.png}
    \caption{Comparison of classification accuracy between training with
        the cross entropy (CE) loss or the MCR$^2$ objective, with
    different ratios of the class labels being randomly corrupted.}
    \label{fig:CIFAR10-performance}
\end{figure}

\begin{figure}[t]
    \begin{subfigure}[t]{0.45\textwidth}
        \centering
        \includegraphics[width=\textwidth]{\toplevelprefix/chapters/chapter4/figs/mcr_pca_individual.pdf}
        \caption{PCAs of features for individual classes learned
        by MCR$^2$.}
        \label{fig:MCR-pca}
    \end{subfigure}
    \hfill
    \begin{subfigure}[t]{0.45\textwidth}
        \centering
        \includegraphics[width=\textwidth]{\toplevelprefix/chapters/chapter4/figs/ce_pca_individual.pdf}
        \caption{PCAs of features for individual classes learned
        by Cross Entropy.}
        \label{fig:CE-pca}
    \end{subfigure}

    \caption{\small Comparison of the principal components of learned
    features from MCR$^2$ versus those from cross entropy.}
    \label{fig:feature-pca}
\end{figure}

\begin{figure*}[t]
    \begin{center}

        \includegraphics[width=0.45\textwidth]{\toplevelprefix/chapters/chapter4/figs/heatmap_mcr2.pdf}
        \hspace{0.25cm}
        \hfill
        \includegraphics[width=0.45\textwidth]{\toplevelprefix/chapters/chapter4/figs/heatmap_ce.png}
        \caption{\small Cosine similarity between learned features by
            using the MCR$^2$ objective  (\textbf{left}) and CE loss
        (\textbf{right}).}
        \label{fig:heatmap-plot}
    \end{center}
    \vskip -0.1in
\end{figure*}

\Cref{fig:feature-pca} shows the distribution of singular values of
the learned features $\vZ_k$ per class. From the distribution of the
singular values shown in \Cref{fig:MCR-pca}, we see that the
features $\vz$ learned from MCR$^2$ for each class are distributed on
a subspace of dimension around 10-12, whereas the features from cross
entropy \Cref{fig:CE-pca} are of much lower dimension.\footnote{If
    one continue to train with more iterations, the CE features will may
    even collapse to a single dimension, known as the neural collapse
phenomenon discussed earlier in \Cref{remark:neural-collapse}. }

\Cref{fig:heatmap-plot} shows the cosine similarities between the
learned features sorted by class. We also compare the similarities of
the learned features by using the cross-entropy (CE)
\eqref{chap4-eqn:cross-entropy} and the MCR$^2$ objective
\eqref{eqn:maximal-rate-reduction}. From the plots, one can clearly
see that the features learned by the  MCR$^2$ objective for different
classes are much more incoherent (or independent) than features
learned by using cross-entropy loss. Within each class, the features
also are less coherent hence more diverse, as indicated by the
distribution of their singular values shown in \Cref{fig:MCR-pca}.

We visualize the images that are mapped along different singular
vectors in the feature subspace of each class, which we call the
``principal'' images. \Cref{fig:pca-images} shows these
principal images along each of the top-10 singular vectors. As we can
see, the images along each singular vector are very similar in their
appearances and structures and images along different singular
vectors are very distinct from each other.
\begin{figure}
    \centering
    \includegraphics[width=\linewidth]{chapters/chapter4/figs/visual_2class.pdf}
    \caption{Top-10 ``principal'' images for  class - ``Bird'' and
    ``Ship'' in the CIFAR-10 image dataset.}
    \label{fig:pca-images}
\end{figure}

Despite the fact that the MCR$^2$ objective seems to allow us to
learn better representations, there has been an apparent lack of
justification of the network architectures used in the above
experiments. It is yet unclear why the network adopted here (the
ResNet-18) is suitable for representing the map $f(\vx, \theta)$, let
alone for interpreting the layer operators and parameters $\theta$
learned inside. In \Cref{ch:deep-networks}, {\em we will
    show how to derive network architectures and components entirely as a
``white box'' from the desired objective (say the rate reduction)}.

\subsection{Unsupervised Representation Learning}
\label{sec:SimDINO}
% \yima{Ziyang, probably introdue the objective of SimDINO here? and
% featuring some experimental justificaiton and leave the
% implementation details to Chapter 7 (or the new Chapter 8).}

As we have seen in the above section, the class information makes
learning a structured representation of the (mixed) data distribution
relatively easy. However, such information is often not available, at
least not at a massive scale.

Now we present how to extend the MCR$^2$ principle to the
unsupervised setting. Contrary to supervised learning where the class
labels are known, in unsupervised learning, the group memberships of
the data samples are unknown. In this case, we can apply certain
class of augmentations or transformations, say $\tau$ with a
distribution $P_{\tau}$, to each sample. For example, for images,
augmentations typically include translations, rotations, or cropping
of each image sample. We may view the augmented samples of each data
point as belonging to the same group or cluster.

More specifically, given a set of $N$ data samples $\vX =
[\vx_1,\dots,\vx_N]$, we generate $m$ augmented views
$\{\tau_i\}_{i=1}^m$ for each data point $\vx_k$ , denoted as
$\tilde{\vX}_{k} = [\tau_1(\vx_k), \tau_2(\vx_k), \dots, \tau_m(\vx_k)]$.
We treat the augmented samples $\tilde{\vX}_{i}$ as the $i$-th class
with $K=N$ groups/clusters in total. Now, we can apply the MCR$^2$
principle as we did before in the supervised learning case.
Specifically, we learn the feature mapping $f(\vx,\theta)$ by
maximizing the coding rate reduction as follows:
\begin{equation}
    \begin{aligned}
        \max_{\theta } & \;  \Delta R_{\epsilon}\big(\vZ \big) \doteq
        R_{\epsilon}(\vZ) - R_{\epsilon}^c(\vZ ), \\
        \mbox{s.t.} & \ \ \, \vZ = f(\tilde{\vX}, \theta),\  \|\vZ_k\|_F^2
        = m,\ k=1,\dots,N.
    \end{aligned}
    \label{eqn:maximal-rate-reduction-unsupervised}
\end{equation}
Here, $\tilde{\vX} = [\tilde{\vX}_1,\dots,\tilde{\vX}_N]$ and $\vZ_k =
f(\tilde{\vX}_k,\theta)$ denotes the features of the augmented samples
of $\vx_k$ for each $k=1,\dots,N$.

Let us again adopt ResNet-18 as the feature mapping $f(\vx,\theta)$
and optimize the network parameters $\theta$ to maximize the
unsupervised objective \eqref{eqn:maximal-rate-reduction-unsupervised}.
After training on the CIFAR-10 dataset, we visualize the evolution of
the two rates and their difference (for both training and test data)
over epochs of training in \Cref{fig:train-test-loss-selfsup-mcr}.
Similar to the supervised learning case, after an initial phase,
$R_\epsilon$ gradually increases while $R^c_\epsilon$ decreases,
indicating that features $\bm Z$ are expanding as a whole while each
cluster $\bm Z_k$ is being compressed. Compared to the supervised
case shown in \Cref{fig:train-test-loss}, we can observe that, in the
unsupervised case, the total coding rate $R_\epsilon$ expands rapidly
and the overall MCR$^2$ objective saturates at a local optimum. To
alleviate this, we can alter the optimization strategy by controlling
the dynamics of $R_\epsilon$ and $R^c_\epsilon$. Specifically, we can
rescale the rates by replacing the coding rate $R_\epsilon(\vZ)$ in
\eqref{eqn:maximal-rate-reduction-unsupervised} with
$$\tilde{R}_\epsilon(\vZ) \doteq \frac{1}{2\gamma_1} \log \det
\left(\I + \frac{\gamma_2 d}{N m \epsilon^{2}}\vZ \vZ ^{\top}\right).$$
By setting $\gamma_1 = \gamma_2 = N$, we arrive at a modified
learning objective, referred to as MCR$^2$-{\scriptsize CTRL}, whose
training dynamics are shown in
\Cref{fig:train-test-loss-selfsup-mcr-ctrl}. As we can see, with this
controlled MCR$^2$ objective, features are first compressed and then
gradually expanded, leading to a higher value of the overall
objective $\Delta R_{\epsilon}(\vZ)$.

To assess the quality of the learned representations and verify
whether they have good subspace structures, we adopt a subspace
clustering algorithm EnSC \cite{you2016oracle}. For evaluation
metrics, we consider normalized mutual information (NMI), clustering
accuracy (ACC), and adjusted rand index (ARI) \footnote{Refer to the
Appendix in \cite{yu2020learning} for their specific definitions.}.
For comparison, we consider methods including JULE
\cite{yang2016joint}, RTM \cite{nina2019decoder}, DEC
\cite{xie2016unsupervised}, DAC \cite{chang2017deep}, and DCCM
\cite{wu2019deep}. The clustering results are summarized in
\Cref{fig:unsupervised-clustering-results}. As we can see, the
features learned by maximizing the MCR$^2$-{\scriptsize CTRL}
objective achieve better clustering accuracy than those learned by
other algorithms.
\begin{figure}[t]
    \begin{subfigure}[t]{0.45\textwidth}
        \centering
        \includegraphics[width=\textwidth]{\toplevelprefix/chapters/chapter4/figs/loss_selfsup_mcr.png}
        \caption{MCR$^2$}
        \label{fig:train-test-loss-selfsup-mcr}
    \end{subfigure}
    \hfill
    \begin{subfigure}[t]{0.45\textwidth}
        \centering
        \includegraphics[width=\textwidth]{\toplevelprefix/chapters/chapter4/figs/loss_selfsup_mcr_ctrl.png}
        \caption{MCR$^2$-ctrl}
        \label{fig:train-test-loss-selfsup-mcr-ctrl}
    \end{subfigure}
    \caption{\small Evolution of the rates of MCR$^2$ in the training
    process for unsupervised learning on CIFAR-10.}
    \label{fig:train-test-loss-selfsup}
\end{figure}

\begin{figure}
    \centering
    \includegraphics[width=\linewidth]{chapters/chapter4/figs/clustering_results.png}
    \caption{Comparison of clustering performance accuracy between
        training with the MCR$^2$-{\scriptsize CTRL} objective and other
        unsupervised learning algorithms on CIFAR-10, CIFAR-100, and
    STL-10 datasets.}
    \label{fig:unsupervised-clustering-results}
\end{figure}

There is one main drawback though: the MCR$^2$-based unsupervised
learning method requires a costly $\log \det$ operation of $\O(d^3)$
complexity \textit{per sample} when computing $R^c_\epsilon$. This
quickly becomes computationally prohibitive as we try to scale to
larger datasets. Instead of minimizing the exact coding rate for each
group of augmented samples $\tilde{\vX}_{i}$, we can approximate the
effect of compressing features in each group by maximizing the {\em
cosine similarity} between the features of the augmented
samples:\footnote{We leave as an exercise for the readers to verify
    under what conditions that minimizing the per class coding rate
    $R^c(\vZ)$ is (approximately) equivalent to maximizing the cosine
similarity among the features $\vZ$.}

\begin{equation}
    \begin{aligned}
        \max_{\theta } & \;  R_{\epsilon}(\vZ) +
        \frac{\gamma}{N}\sum_{k=1}^N \sum_{i \neq j}
        \vz_{k,j}^\top \vz_{k,i}, \\
        \mbox{s.t.} &  \ \ \, \vz_{i,j} = f(\tau_j(\vx_i),\theta),\
        \|\vz_{i,j}\|_2^2                                          = 1,
    \end{aligned}
    \label{eqn:maximal-rate-reduction-unsupervised-simdino}
\end{equation}
where $\gamma > 0$ is a hyperparameter that controls the strength of
the cosine similarity term, as it differs from the coding rate in
magnitude. This objective avoids the costly $\log \det$ operation and
is much more scalable. In fact, this objective is closely related to
popular (unsupervised) {\em contrastive learning} methods such as
SimCLR \cite{chen2020simple} and DINO \cite{caron2021emerging}, which
have been engineered to learn high-quality representations from
massive amount of images.

As we will show later in
\Cref{sub:contrastive_learning_objective} of
\Cref{ch:applications}, this objective
\eqref{eqn:maximal-rate-reduction-unsupervised-simdino} can be
implemented very efficiently to learn a highly structured
representation $\vZ$ that is highly linear and discriminative.
% \yima{Robin, maybe some experimental results that illustrate
% qualitatively the learned features are discriminative and linear?
% Say some results on linear probing or visual results on
% segmentation? (it is ok to say leaving the implementation details
% to Chapter 8.) This will help justify RAE later in Chapter 6.}
The learned feature representation leads to state-of-the-art
performance on large real-world datasets comparable or even better
than features learned by DINO, but without the heavy engineering
efforts needed for DINO. Hence, we refer to this objective as {\em
Simplified DINO} (SimDINO). As a preview,
\Cref{tab:dino_imagenet_linear_probing_teaser} show that SimDINO
outperforms DINO on ImageNet-1k classification using both
\(k\)-nearest neighbor accuracy and linear probing (i.e. training a
linear classifier on pretrained models). This demonstrates that the
objective leads to highly structured and linearized representations.
We can also utilize these representations for downstream tasks such
as unsupervised segmentation\footnote{Specifically, we adopt MaskCut
    \cite{wang2023cut}, an effective unsupervised approach of extracting
    features from a frozen vision backbone for object detection and
instance segmentation.}, as illustrated in
\Cref{fig:maskcut_segmentation}. More implementation details and
experimental results of SimDINO will be presented in
\Cref{ch:applications}.

\begin{table}
    \centering
    \begin{tabular}{@{}lcccc@{}} %
        \toprule
        Method & Model & Epochs & 20-NN & Linear Probing
        \\
        \midrule
        DINO & ViT-B & 100 & 72.9 & 76.3 \\
        SimDINO & ViT-B & 100 & \bf 74.9 & \bf 77.3 \\
        DINO & ViT-L & 100 & -- & -- \\
        SimDINO & ViT-L & 100 & \bf 75.6 & \bf 77.4 \\

        \bottomrule
    \end{tabular}\smallskip
    \caption{\small\textbf{Classification performance} on ImageNet-1k
        test set for DINO and SimDINO, using both \(k\)-nearest neighbor
        accuracy (\(k = 20\)) and linear probing. SimDINO
        outperforms DINO in
        both accuracy and training stability. DINO training on ViT-L is
    unstable and quickly diverges, while SimDINO trains stably.}
    \label{tab:dino_imagenet_linear_probing_teaser}
\end{table}

\begin{figure}
    \centering
    \includegraphics[width=\linewidth]{chapters/chapter4/figs/maskcut_demo.pdf}
    \caption{\small \textbf{Visualization of unsupervised segmentation
        results} from DINO ViT-B/16 (row 1), SimDINO ViT-B/16 (row 2) and
    SimDINO ViT-L/16 (row 3).}
    \label{fig:maskcut_segmentation}
\end{figure}

\section{Summary and Notes}

\paragraph{Key messages.}
In this Chapter, we have studied basic concepts and ideas behind how
to learn a distribution from finite samples and obtain a {\em
computable} encoding and decoding scheme for the distribution. For
distributions that are continuous and possibly low-dimensional, we
have argued that such a scheme is necessarily {\em lossy} -- hence
the importance of the concept of rate distortion. The noise
introduced through distortion, due to quantization, plays a related
but different role from the noise we have studied in the preceding
chapter. One may think of the noise considered in
\Cref{ch:denoising} as ``outside'' of the submanifold of the support
of the distribution whereas the noise considered in this chapter is
``inside'' the support. Informally, we may refer to the
diffusion/denoising process outside of the support as
``continuation'' and the diffusion/denoising process inside the
support as ``percolation.'' A more unified analysis and understanding
of their compounded effect is due for future study.

In this chapter, we introduced for the first time the concept of
learning a {\em representation} for a distribution. We have
considered arguably the simplest case, the case of classification. If
one views the classification label $\y$ as a fixed simple (grossly
lossy) code for the data $\vx$, we may seek a consistent\footnote{That
    is, we can faithfully recover $\vx$ from $\vz$ as we will make the
    notion of ``consistent'' representations clear later in
\Cref{ch:consistent-representations}.} representation $\vz = f(\vx)$ of
the data $\vx$
such that it tends to maximize the mutual information between the
representation and this simple code:
\begin{equation}
    \max I(\vz, \y) = H(\vz) - H(\vz \mid \y).
\end{equation}
In addition, we require such a representation $\vz$ is not only
consistent to $\vx$ but also much more {\em structured} than the
native data distribution of $\vx$ -- classes are linearly
discriminative. This would make the learned representation much more
efficient to use than the native data distribution for subsequent
tasks such as recognition, segmentation and completion. The maximal
coding rate reduction objective introduced is a natural and
computable realization of these requirements. In
\Cref{ch:conditional-inference}, we will see how to generalize this
objective to a much broader context when $\y$ could be a more general
lossy code (say image caption) or a random variable related to $\vx$.

\section{Exercises and Extensions}

\begin{exercise}
    Please show the following properties of the  $\log\det(\cdot)$ function.
    \begin{enumerate}
        \item Show that
            \begin{align*}
                f(\vX) = \log\det\left(\vX\right)
            \end{align*}
            is a concave function. ({\bf Hint:} The function $f(\vx)$ is
                concave if and only if the one-dimensional function $f(\vx+t\vh)$
            is concave for all $\vx$ and $\vh$.)

        \item Show that
            \begin{align*}
                \log\det(\vI + \vX^\top\vX) = \log\det(\vI + \vX\vX^\top)
            \end{align*}

        \item Let $\vA \in \R^{n\times n}$ be a positive definite matrix.
            Please show that:
            \begin{align}
                \log\det\left(\vA \right) = \sum_{i=1}^n \log(\lambda_i),
            \end{align}
            where $\lambda_1,\lambda_2,\dots,\lambda_n$ are the
            eigenvalues of $\vA$.

        \item  Show that the gradient of the $\log\det(\cdot)$
            function is given by
            \begin{align*}
                \nabla \log\det(\vA) = \vA^{-1},
            \end{align*}
            for a positive definite matrix $\vA$.

    \end{enumerate}

\end{exercise}

% \begin{exercise}
%     Let $\bm Z \in \R^{d\times N}$ and define
%     \begin{equation}
%     R_{\epsilon}(\vZ) := \frac{1}{2}\log\det\left(\I +
%     \alpha\vZ\vZ^{\top}\right),
%     \end{equation}
%     where $\alpha > 0$.
%     \begin{enumerate}
%         \item Please compute the gradient $\nabla R_{\epsilon}(\bm
% Z)$ of the above function.

%         \item Let $\bm D \in \mathbb{R}^{d \times n}$ be arbitrary.
% Please compute the bilinear form of the Hessian of the above function, i.e.,
%         \begin{align*}
%             \nabla^2 R_{\epsilon}(\bm Z)[\bm D, \bm D] :=
% \left\langle \bm D, \lim_{t \to 0} \frac{\nabla R_{\epsilon}(\bm Z
% + t\bm D) - \nabla R_{\epsilon}(\bm Z)}{t}\right\rangle.
%         \end{align*}
%     \end{enumerate}
% \end{exercise}

\begin{exercise}
    Let $\bm Z = [\bm Z_1, \dots, \bm Z_K] \in \mathbb{R}^{d \times N}$,
    where each $\bm Z_k \in \mathbb{R}^{d \times N_k}$ and
    $N = \sum_{k=1}^K N_k$.

    \begin{enumerate}
        \item Please show that
            \begin{align*}
                \log\det\left( \bm I + \alpha \bm Z \bm Z^\top \right)
                \le
                \sum_{k=1}^K
                \log\det\left( \bm I + \alpha \bm Z_k \bm Z_k^\top \right),
            \end{align*}
            where $\alpha > 0$.

        \item Please derive the necessary and sufficient condition
            under which the equality holds.
    \end{enumerate}

\end{exercise}

\begin{exercise}
    Let $\vz_1$ and $\vz_2$ be two vectors of unit length. Let
    $\vZ = [\vz_1,
    \vz_2]$ Show that the following three objectives are equivalent:
    \begin{itemize}
        \item $\max \vz_1^\top \vz_2$.
        \item $\min \|\vz_1 - \vz_2\|_2^2$.
        \item $\max \log \det (\I + \vZ\vZ^\top)$.
    \end{itemize}
    However, they may no longer be exactly equivalent when there are more
    than two vectors in a high-dimensional space. The $\log\det$ is more
    general as it measures the total volume spanned by all the vectors in
    the high-dimensional space.
\end{exercise}

% {\color{blue}
% \begin{exercise}\label{exer:gaussian_rate_distortion}
%     Please show that the rate distortion function for a $X \sim
%     \mathcal{N}(0,\sigma^2)$ source is
%     \begin{align*}
%         \mathcal{R}_{\epsilon}(X) =
%         \begin{cases}
%             \frac{1}{2}\log\left( \frac{\sigma^2}{\epsilon^2}
%             \right),\ & 0 \le \epsilon \le \sigma,\\
%             0,\ & \epsilon > \sigma.
%         \end{cases}
%     \end{align*}
%     {\bf Hint:}  This proof consists of two steps. First, we find a
%     lower bound for the rate distortion function, i.e., for all
%     $\mathbb{E}[(X-\hat{X})^2] \le \epsilon^2$,
%     \begin{align*}
%         I(X;\hat{X}) \ge  \frac{1}{2}\log\left(
%         \frac{\sigma^2}{\epsilon^2} \right).
%     \end{align*}
%     Second, we construct $p(\hat{X} \mid X)$ to prove this lower
%     bound is achievable. If $\epsilon \le \sigma$, we choose
%     \begin{align*}
%         X = \hat{X} + Z,\ \hat{X} \sim \mathcal{N}(0,\sigma^2 -
%         \epsilon^2),\ Z \sim \mathcal{N}(0,\epsilon^2),
%     \end{align*}
%     where $\hat{X}$ and $Z$ are independent. If $\epsilon > \sigma$,
%     we choose $\hat{X} = 0$ with probability 1.
% \end{exercise}

\begin{exercise}\label{exer:KKT}
    Please derive the rate-distortion function for
    the Gaussian random vector introduced in
    \Cref{ex:gaussian_rate_distortion}.
    % \textbf{Hints:} The proofs consists of the following steps:
    % \begin{enumerate}
    %     \item We first find a lower bound for the rate distortion
    %         function, i.e., for all $\mathbb{E}[\|\bm{x}-\hat{\bm
    %         x}\|_2^2] \le \epsilon^2$,
    %         \begin{align}
    %             I(\bm x; \hat{\bm x}) \ge \sum_{i=1}^d
    %             \max\left\{\frac{1}{2}\log\left(
    %             \frac{\sigma_i}{\lambda_i} \right), 0\right\},
    %         \end{align}
    %         where $\lambda_i = \mathbb{E}[(x_i - \hat{x}_i)^2]$.

    %     \item Second, we reduce the problem of finding the rate
    %         distortion function to the followign problem:
    %         \begin{align}
    %             \min_{\lambda_1,\dots,\lambda_d}  \sum_{i=1}^d
    %             \max\left\{\frac{1}{2}\log\left(
    %             \frac{\sigma_i}{\lambda_i} \right), 0 \right\}\quad
    %             \mathrm{s.t.}\ \sum_{i=1}^d \lambda_i = \epsilon^2.
    %         \end{align}

    %         % \item Consider setting up the Lagrangian
    %         % \[
    %         %     \mathcal{L}(\lambda_1,\dots,\lambda_d,\nu) =
    %         % \frac12 \sum_{i=1}^d \log \frac{\sigma_i^2}{\lambda_i}
    %         %     + \nu \left(\sum_{i=1}^d \lambda_i -
    %         % \epsilon^2\right),
    %         % \]
    %         % where \(\lambda_i\) are the variances of the
    %         % reconstruction errors along each principal direction.

    %         % \item Compute the partial derivative \(\partial
    %         % \mathcal{L}/\partial \lambda_i = 0\) to find a relation
    %         % between \(\lambda_i\) and the Lagrange
    % multiplier \(\nu\).

    %         % \item Apply the constraints \(0 \le \lambda_i \le
    %         % \sigma_i^2\) to identify which directions are saturated
    %         % and which follow the water-filling allocation.

    %         % \item Solve for the ``water level'' \(\lambda\) using
    %         % \(\sum_{i=1}^d \lambda_i = \epsilon^2\), and verify
    %         % that \(\lambda_i = \min(\sigma_i^2, \lambda)\).
    % \end{enumerate}
\end{exercise}
% }

\end{document}
